\documentclass[11pt]{article}
\usepackage[margin=1.1in]{geometry}
\usepackage{amsmath,amssymb,amsthm,mathtools,enumitem}
\usepackage{xcolor}
\usepackage{xr-hyper}
\usepackage[colorlinks=true,linkcolor=blue!50!black,citecolor=blue!50!black]{hyperref}

\newtheorem{theorem}{Theorem}[section]
\newtheorem{lemma}[theorem]{Lemma}
\newtheorem{proposition}[theorem]{Proposition}
\newtheorem{corollary}[theorem]{Corollary}
\theoremstyle{definition}
\newtheorem{definition}[theorem]{Definition}
\newtheorem{conjecture}[theorem]{Conjecture}
\newtheorem{claim}[theorem]{Claim}
\newtheorem{assumption}[theorem]{Assumption}
\theoremstyle{remark}
\newtheorem{remark}[theorem]{Remark}
\newtheorem{question}[theorem]{Question}

\newcommand{\cA}{\mathcal A}
\newcommand{\cB}{\mathcal B}
\newcommand{\cC}{\mathcal C}
\newcommand{\Om}{\Omega}
\newcommand{\Lam}{\Lambda}
\newcommand{\ph}{\varphi}
\newcommand{\bfB}{\mathbf B}
\newcommand{\bfA}{\mathbf A}
\newcommand{\bfC}{\mathbf C}
\newcommand{\bfE}{\mathbf E}
\newcommand{\fus}{\boxtimes}
\newcommand{\cfus}{\circledast}
\newcommand{\ip}[2]{\langle #1,#2\rangle}
\newcommand{\Id}{\mathrm{Id}}
\newcommand{\Wk}{\mathbf 1^{\mathrm w}_{\cA}}
\newcommand{\Wkt}[1]{\mathbf 1^{\mathrm w}_{\cA,#1}}

\title{Defect fusion, boundaries and junctions in conformal QFT\\[2pt]\large Applications of the fusion theorem for $Q$-system defects}
\author{Notes prepared for A.~Bailey}
\date{18 October 2026}
\externaldocument[N-]{connes_fusion_notes}
\newcommand{\nref}[1]{N\,\ref{N-#1}}


\begin{document}
\maketitle
\begin{abstract}
\noindent This companion applies the results of the main notes (\emph{Fusion of $Q$-system defects and juxtaposition of boundary conditions}, cited as N) to two-dimensional conformal quantum field theory. It relates the defects $D_{\mathbf m}$ to the boundary nets of Bischoff--Kawahigashi--Longo--Rehren by time-zero restriction, proves that fusion through an intermediate phase gives the summand $\mathbf m_1\otimes_{\bfB}\mathbf m_2$, derives the defect--boundary fusion rule for hard boundaries, and studies junctions, crossings of transparent boundaries, corners, lightlike boundaries and interfaces, and hard boundaries meeting at a point.
\end{abstract}
\noindent\textbf{Conventions.} A reference such as ``Theorem~N\,9.4'' is to the numbering of the main notes (file \texttt{connes\_fusion\_notes}); unmarked references are to this companion. Notation is that of the main notes. All proofs are by hand and have not been machine-checked; status remarks mark what is proved, quoted or conjectural.

\tableofcontents

\section{Introduction}\label{sec:apps}
This section connects the defects $D_{\mathbf m}$ to the two-dimensional boundary nets of [BKLR2] and to boundary CFT. Statements marked \emph{dictionary} depend on identifications with constructions in the literature that are recalled, not re-proved, here.

\section{Time-zero restriction}\label{sec:tz}
Let $A$ be a chiral net on the light ray and $\cA_2(O)=A(I_+)\otimes A(I_-)$ on double cones $O=I_+\times I_-$ of two-dimensional Minkowski space (light-cone coordinates $x_\pm=t\pm x$). A boundary condition in the sense of [BKLR2, Def.~3.2, Eq.~(3.4)] is a pair of local irreducible extensions $\cB^L,\cB^R\supset\cA_2$ on a common Hilbert space, with $D(O)=\cB^L(O)\vee\cB^R(O)$ left-local with respect to $\cB^R$ and right-local with respect to $\cB^L$; the physical theory uses $\cB^L$ to the left and $\cB^R$ to the right of the timelike line $x=0$. The universal construction [BKLR2, Prop.~5.1] is the extension of $\cA_2$ by the braided product $Z[\bfA]\times Z[\bfB]$ (for $\cB^L=\cA_2\rtimes Z[\bfA]$, $\cB^R=\cA_2\rtimes Z[\bfB]$), and the irreducible boundary conditions are its minimal central projections $E_{\mathbf m}$ [BKLR2, Cor.~5.3, Thm.~5.5].

\begin{lemma}[time zero]\label{lem:tz}
For an interval $J$ of the time-zero line let $O_J$ be its causal completion. Then $\cA_2(O_J)=A(J)\otimes A(rJ)=\cA(J)$ with $\bar A(J)=A(rJ)$ as in Lemma~\nref{lem:conjnet}. Localized endomorphisms $\rho\otimes\sigma$ of $\cA_2$ localized in $O_J$ are exactly the localized endomorphisms of $\cA$ localized in $J$, so $\mathrm{DHR}(\cA_2)$ and $\mathrm{DHR}(\cA)\supset\mathcal C\boxtimes\mathcal C^{\mathrm{opp}}$ have the same $Q$-systems, and for a $Q$-system $Q$ localized in $J_0$, $(\cA_2\rtimes Q)(O_J)=\langle\cA(J),v\rangle=(\cA\rtimes Q)(J)$ for $J\supset J_0$.
\end{lemma}
\begin{proof}
At $t=0$, $x_+=x$ and $x_-=-x$, so $O_J=J\times(-J)$. The extension algebra on a double cone is generated by the observables and one charged intertwiner localized in it, by transportability; both sides are defined by the same formula.
\end{proof}
Write $\cA_{\bfA}:=\cA\rtimes Z[\bfA]$ for the time-zero net of the phase $\bfA$, and $D^{\rm tz}(J):=\cA_{\bfA}(J)\vee\cA_{\bfB}(J)=\langle\cA(J),v^{\bfA}v^{\bfB}\rangle$ for the time-zero universal construction ($J$ containing both localization regions).

\begin{proposition}[pinning]\label{prop:pinning}
Let $Q''=Z[\bfA]^h\times Z[\bfB]$ as in Theorem~\nref{thm:net}, localized in the strip $[0,1]$ with $Z[\bfA]^h$ to the left of $Z[\bfB]$. For every genuinely bicolored $I$, $D_{Q''}(I)=D^{\rm tz}(I^\sharp)$, and the central projections $E_{\mathbf m}$ of Theorem~\nref{thm:fullc} are those of [BKLR2]. Hence $D_{\mathbf m}(I)=E_{\mathbf m}D^{\rm tz}(I^\sharp)$ is the time-zero algebra of the irreducible boundary condition $\mathbf m$. The defect $D_{\mathbf m}$ differs from the physical boundary net $P_{\mathbf m}$ (which assigns $\cA_{\bfA}(J)$, resp.\ $\cA_{\bfB}(J)$, to intervals $J$ left, resp.\ right, of the boundary point) only on monochrome intervals, where $D_{\mathbf m}$ assigns $\cA(J)$: $D_{\mathbf m}$ is $P_{\mathbf m}$ \emph{pinned} to $\cA$ away from the boundary.
\end{proposition}
\begin{proof}
Both algebras are generated by $\cA(I^\sharp)$ and $v''=\alpha_h(v^{\bfA})v^{\bfB}$; by Lemma~\ref{lem:tz} and transportability $\cA_{\bfA}(I^\sharp)=\langle\cA(I^\sharp),\alpha_h(v^{\bfA})\rangle$ and similarly for $\bfB$, and $\alpha_h(v^{\bfA})=w^{\bfB*}v''$, $v^{\bfB}=\alpha_h(w^{\bfA})^*v''$ (Theorem~\nref{thm:net}). The extension is the same $Q$-system, so its centre is the same. Left-locality of $\cA_{\bfA}$ with respect to $\cA_{\bfB}$ is the commutation relation $v^{\bfB}v^{\bfA}=\iota(\varepsilon^\pm)v^{\bfA}v^{\bfB}$ that defines the braided product; which sign arises does not matter (Remark~\nref{rem:signs}).
\end{proof}
So Theorem~\nref{thm:A} says that \emph{BDH fusion of the pinned phase strips $D_{Z[\bfA]}$, $D_{Z[\bfB]}$ is BKLR's universal construction.}

\section{Fusion through the intermediate phase}
For the unpinned boundary nets, juxtaposition of $\mathbf m_1\colon\bfA\to\bfB$ and $\mathbf m_2\colon\bfB\to\bfC$ means BDH fusion over the net $\cA_{\bfB}$ of the intermediate phase. We use one more standard fact:
\begin{itemize}[leftmargin=2.6em]
\item[(L4)] For a finite-index inclusion $C\subset B$ with expectation and Pimsner--Popa basis $\{b_i\}$, and a right $B$-module $H$, left $B$-module $K$, $H\fus_BK$ is the range of the projection $q=\sum_i\rho_H(b_i)\fus\lambda_K(b_i^*)$ on $H\fus_CK$ (relative tensor product over a finite-index extension as a compression).
\end{itemize}

\begin{theorem}[fusion through the phase $\bfB$]\label{thm:through}
Let $P_{\mathbf m_1}$, $P_{\mathbf m_2}$ be the unpinned time-zero boundary nets, with vacuum sectors $H_1=H_0(D_{\mathbf m_1})$, $H_2=H_0(D_{\mathbf m_2})$ (the algebras on genuinely bicolored intervals agree, Proposition~\ref{prop:pinning}). Then, on genuinely bicolored intervals,
\[
(P_{\mathbf m_1}\cfus_{\cA_{\bfB}}P_{\mathbf m_2})(I)\ \cong\ D_{(\mathbf m_1,\mathbf m_2)}(I),
\]
the summand $\mathbf m'=\bfB$ of Theorem~\nref{thm:main}. In particular its relative commutant is $\mathrm{End}_{\bfA\text{-}\bfC}(\mathbf m_1\otimes_{\bfB}\mathbf m_2)$, which is (E2) in the literal sense of [BKLR2].
\end{theorem}
\begin{proof}
(a) By (L4) with $C=\cA(J)\subset B=\cA_{\bfB}(J)$, $H_1\fus_{\cA_{\bfB}(J)}H_2=q(H_1\fus_{\cA(J)}H_2)$.
(b) The generators of the fused algebra over $\cA$, $D(K)=\langle\cA(K),v''_1\rangle$ and $E(L)=\langle\cA(L),v''_2\rangle$, commute with the right action of $\cA_{\bfB}(\tilde J)$ on $H_1$ and the left action of $\cA_{\bfB}(J)$ on $H_2$ respectively: the $\bfB$-charged fields commute among themselves because $\cA_{\bfB}$ is local ($Z[\bfB]$ is commutative), and with the $\bfA$-, resp.\ $\bfC$-charged fields by left-, resp.\ right-locality. Hence $D(K)=P_{\mathbf m_1}(I^{++})\cap\cA_{\bfB}(\tilde J)'$ (Haag duality over $\cA_{\bfB}$ gives the same algebra as over $\cA$), likewise for $E(L)$, and $q$ commutes with $F=D(K)\vee E(L)$. So the fused algebra over $\cA_{\bfB}$ is the compression $qF$.
(c) In $Q$-system form (L4) reads: for modules $H,K$ over the special Frobenius algebra $Z[\bfB]$, $H\otimes_{Z[\bfB]}K$ is the range of $d^{-1}(r_H\otimes l_K)(1\otimes x\eta\otimes1)$, the action of the normalized Casimir $c$. Under Theorem~\nref{thm:net} and Corollary~\nref{cor:complete}, $F\cong pD_{\Theta_4}$, the right $\cA_{\bfB}$-action on $H_1$ and the left one on $H_2$ become multiplication by the second and third copies of $Z[\bfB]$ in $\Theta_4$, and $q$ becomes multiplication by $\eta\otimes c\otimes\eta$. By Lemma~\nref{lem:casimir}, $c=E_{\bfB}$, so $q=q_{E_{\bfB}}$, and $qF\cong pq_{E_{\bfB}}D_{\Theta_4}\cong D_{(\mathbf m_1,\mathbf m_2)}$ by the proof of Theorem~\nref{thm:E2}.
\end{proof}
\emph{Status.} (a) is a standard fact; (b) is proved; (c) is Lemma~\nref{lem:casimir} together with the identification of the two middle $\cA_{\bfB}$-actions with the middle factors of $\Theta_4$, which is how Theorem~\nref{thm:net} glues the strips. That the fused object over $\cA_{\bfB}$ satisfies all defect axioms is not addressed. Physically: fusing over the ambient chiral net $\cA$ lets the shrunken strip of phase $\bfB$ carry any $\bfB$--$\bfB$ boundary condition $\mathbf m'$ (Theorem~\nref{thm:main}(iii)); fusing through the phase $\bfB$ keeps exactly $\mathbf m'=\bfB$.

\section{Hard boundaries}\label{sec:hard}
A hard boundary of a phase $\bfB$ is a timelike line beyond which there is nothing, at which only one chiral stress tensor survives. Longo--Rehren [LR04] construct the boundary net on the half-plane from a chiral extension $A\subset\cB_Y$ given by a $Q$-system $Y$ in $\mathcal C$; the bulk is the full centre $Z[Y]$ [LR09, BKL16]. Unfolding the half-plane to the light ray (time-zero points $x>0$ go to the pair $\pm x$, the boundary to $0$), the algebra of an interval $[0,b)$ touching the boundary is $\cB_Y((-b,b))=\langle A((-b,b)),v_Y\rangle$, the chiral defect $D^{\rm ch}_Y$ at the colour point $0$ (a $Q$-system defect of the chiral net $A$, to which Theorem~\nref{thm:A} also applies).

\emph{Input (H) (quoted, used only for interpretation).} [LR04, LR09, BKL16; FRS for the categorical form] The LR boundary net with chiral algebra $Y$ has bulk $Z[Y]$; the irreducible hard boundary conditions of the phase $\bfB$ are the simple left $\bfB$-modules $\mathbf N$, with chiral algebra $Q_{\mathbf N}:=\underline{\mathrm{End}}(\mathbf N)=\bar{\mathbf N}\otimes_{\bfB}\mathbf N$ (Morita equivalent to $\bfB$, so $Z[Q_{\mathbf N}]=Z[\bfB]$). A reducible chiral algebra $\underline{\mathrm{End}}(X)$, $X=\bigoplus_jK_j\mathbf M_j$, is the boundary condition $\bigoplus K_j\mathbf M_j$.

\subsubsection*{Unfolding}
\begin{lemma}[unfolding]\label{lem:unfold}
For an interval $I\subset(0,\infty)$, $\cA(I)=A(I)\otimes A(-I)\cong A(I)\vee A(-I)\subset A(\mathrm{conv}(I\cup-I))$ (split property). The assignment $U(\rho\otimes\bar\sigma):=\rho\circ\sigma^{(-I)}$, where $\sigma^{(-I)}$ is $\sigma$ localized in $-I$, and $U(T\otimes S):=TS$ on intertwiners, is a faithful strict $*$-tensor functor from the endomorphisms of $\cA$ localized in $I$ to the endomorphisms of $A$ localized in $\mathrm{conv}(I\cup-I)$, equivalent to $F_{\mathbf 1}\colon\mathcal Z\to\mathcal C$. For a $Q$-system $P$ localized in $I$, $D_P$ restricted to intervals in $(0,\infty)$ is the two-interval subalgebra $\langle A(I'^\sharp)\vee A(-I'^\sharp),v\rangle$ of the chiral extension $A\rtimes U(P)$.
\end{lemma}
\begin{proof}
$\rho$ is localized in $I$ and $\sigma^{(-I)}$ in $-I\subset I'$, so they commute and $U(\theta_1)U(\theta_2)=\rho_1\sigma_1\rho_2\sigma_2=\rho_1\rho_2\sigma_1\sigma_2=U(\theta_1\theta_2)$. For $T\in\mathrm{Hom}(\rho,\rho')\subset A(I)$, $S\in\mathrm{Hom}(\sigma,\sigma')\subset A(-I)$ the chiral tensor product is $T\rho(S)=TS$ since $\rho$ is trivial on $A(-I)$. Restricted to $A(I)\vee A(-I)$, $\rho\sigma^{(-I)}$ acts as $\rho\otimes\sigma$. The identification with $F_{\mathbf 1}$ ($\rho\boxtimes\sigma\mapsto\rho\sigma$) is the standard one. The last statement is the definition of $D_P$ transported by $U$.
\end{proof}

\subsubsection*{The Cardy cap}
\begin{lemma}[the trivial boundary absorbs a Cardy strip]\label{lem:cardy}
Let $E=(-b,-a)\cup(a,b)$, $L=(-b,b)$, $K=(-a,a)$. The two-interval inclusion $A(E)\subset\hat A(E):=A(K)'\cap A(L)$ is the extension of $\cA((a,b))$ by the $Q$-system $Z[\mathbf 1]=\bigoplus_ii\otimes\bar i$ [KLM]; this is the bulk of the Cardy phase next to its trivial boundary (LR boundary net with $Y=1$). Its charged generator, realized in the chiral extension $A\rtimes U(Z[\mathbf 1])=A\rtimes F_{\mathbf 1}I_{\mathbf 1}(1)$, is mapped into $A(L)$ by the counit $\varepsilon\colon F_{\mathbf 1}I_{\mathbf 1}(1)=\bigoplus_i\bar i\otimes i\to1$; that is, the Cardy bulk $\hat A(E)$ is the image of the pinned Cardy strip $D_{Z[\mathbf 1]}$ under $\varepsilon$.
\end{lemma}
\begin{proof}
Put $\theta=\bigoplus_i\bar i^{(-I)}i^{(I)}$, $I=(a,b)$, the chiral endomorphism $U(\bigoplus_ii\otimes\bar i)$, localized in $E$. By [KLM], $\hat A(E)=\langle A(E),v\rangle$ for an isometry $v\in\hat A(E)\subset A(L)$ with $va=\theta(a)v$ for $a\in A(E)$, the relations of the $Q$-system $Z[\mathbf 1]$ holding. Since $v\in A(K)'$ and $\theta$ is trivial on $A(K)$, $va=\theta(a)v$ also for $a\in A(K)$, hence for all $a\in A(L)=A(E)\vee A(K)$ (strong additivity). So $v$ defines a homomorphism $A(L)\rtimes F_{\mathbf 1}Z[\mathbf 1]\to A(L)$ which is the identity on $A(L)$. Homomorphisms of extensions that fix the base correspond to $Q$-system morphisms [BKLR, \S4], here an algebra map $\chi\colon F_{\mathbf 1}I_{\mathbf 1}(1)\to1$. Under the monoidal adjunction $F_{\mathbf 1}\dashv I_{\mathbf 1}$, algebra maps $F_{\mathbf 1}(Z)\to X$ correspond to algebra maps $Z\to I_{\mathbf 1}(X)$, so $\chi$ corresponds to a unital algebra endomorphism $\alpha$ of $Z[\mathbf 1]$, and $\varepsilon$ to the identity. $Z[\mathbf 1]$ is haploid, hence simple as an algebra (its bimodule endomorphisms are central elements of $\mathrm{Hom}(1,Z[\mathbf 1])=\mathbb C$), so $\alpha$ is injective, hence an automorphism. Thus $\chi=\varepsilon\circ F_{\mathbf 1}(\alpha)$. Precomposing with an automorphism changes neither the image nor the kernel, so the Cardy bulk is the image of the pinned strip under $\varepsilon$, up to relabelling by $\alpha$. ($\alpha$ amounts to a choice of phases of the transporters, i.e.\ a tensor automorphism of $\mathrm{Id}_{\mathcal C}$ [Davydov].)
\end{proof}
This is the boundary counterpart of Lemma~\nref{lem:unit}: there, the multiplication $Z[\bfB]\times Z[\bfB]\to Z[\bfB]$ glues a pinned $\bfB$-strip into the $\bfB$-phase; here, $\varepsilon$ glues a pinned Cardy strip into the Cardy bulk at the trivial boundary.

\subsubsection*{Capping a stack and defect--boundary fusion}
A stack of transparent defects ending in the Cardy phase, placed in front of the trivial boundary, is described by a pinned $\cA$-defect $D_P$ with $P=W\times Z[\mathbf 1]$, $W$ the rest of the stack. We define its boundary-touching algebra on $[0,b)$ as the algebra generated by the chiral boundary algebra $A((-b,b))$ and $U(v_P)$, i.e.\ the chiral extension $A\rtimes U(P)$, after gluing every pinned strip into the adjacent phase (through-phase fusion, Theorem~\ref{thm:through}, and Lemma~\ref{lem:cardy} for the Cardy strip).

\begin{proposition}[capping]\label{prop:cap}
$U(W\times Z[\mathbf 1])=F_{\mathbf 1}(W)\otimes_{\rm hb}F_{\mathbf 1}I_{\mathbf 1}(1)$ (product via the half-braiding), and $1\otimes\varepsilon\colon F_{\mathbf 1}(W)\otimes_{\rm hb}F_{\mathbf 1}I_{\mathbf 1}(1)\to F_{\mathbf 1}(W)$ is a surjective algebra map. Hence the capped chiral algebra of the stack is $F_{\mathbf 1}(W)$. Without the Cardy gluing (fusion over $\cA$) one gets instead $F_{\mathbf 1}I_{\mathbf 1}(F_{\mathbf 1}W)=\bigoplus_{i}\bar i\otimes F_{\mathbf 1}(W)\otimes i$.
\end{proposition}
\begin{proof}
$U\simeq F_{\mathbf 1}$ is monoidal (Lemma~\ref{lem:unfold}), and the braided product becomes the half-braided product. $\varepsilon$ is an algebra map (counit of a monoidal adjunction), and $1\otimes\varepsilon$ respects the half-braided product by naturality of the half-braiding with respect to $\varepsilon$ and $\widehat X_1=\mathrm{id}$, exactly as in the proof of Lemma~\nref{lem:projalg}. It is split surjective via $1\otimes\eta$. The last formula is the projection formula $W\times I_{\mathbf 1}(1)\cong I_{\mathbf 1}(F_{\mathbf 1}W)$ followed by Proposition~\nref{prop:T8} with $\mathbf X=\mathbf Y=\mathbf 1$.
\end{proof}

\begin{theorem}[defect--boundary fusion]\label{thm:bcft}
Let $\mathbf m\in\mathrm{Irr}({}_{\bfA}\mathcal C_{\bfB})$ be a transparent defect between the phases $\bfA$ and $\bfB$ and $\mathbf N$ a simple left $\bfB$-module. The configuration ``phase $\bfA$ $|$ $\mathbf m$ $|$ phase $\bfB$ $|$ $\mathbf N$ $|$ Cardy phase $|$ trivial boundary'', with all strips glued through their phases, has capped chiral algebra
\[
Y_{(\mathbf m,\mathbf N)}=\underline{\mathrm{End}}_{{}_{\bfA}\mathcal C}(\mathbf m\otimes_{\bfB}\mathbf N),\qquad\mathrm{Hom}(1,Y_{(\mathbf m,\mathbf N)})=\mathrm{End}_{\bfA}(\mathbf m\otimes_{\bfB}\mathbf N).
\]
By (H) this is the hard boundary of the phase $\bfA$ with boundary condition $\mathbf m\otimes_{\bfB}\mathbf N=\bigoplus_jK_j\mathbf M_j$ and zero-mode algebra $\bigoplus_jM_{K_j}(\mathbb C)$. The case $\bfA=\bfB$, $\mathbf m=\bfB$ gives back the boundary $\mathbf N$ of $\bfB$ (consistency with (H)). Fusing over $\cA$ instead of through the Cardy phase gives the reducible boundary condition $\bigoplus_i(\mathbf m\otimes_{\bfB}\mathbf N)\otimes i$.
\end{theorem}
\begin{proof}
By Theorems~\nref{thm:main} and \ref{thm:through} with $\bfC=\mathbf 1$, the configuration through the phase $\bfB$ is the summand $f_{\mathbf m\mathbf N}$ of $W\times Z[\mathbf 1]$, $W=Z[\bfA]\times Z[\bfB]$. By Proposition~\ref{prop:cap} the capped chiral algebra is the corresponding summand of $F_{\mathbf 1}(W)$. By (T1), (T2) and Lemma~\nref{lem:projalg}, $W\cong I_{\bfB}(Y_1)$, $Y_1=\bigoplus_{\mathbf m}\underline{\mathrm{End}}(\mathbf m)$, and by Proposition~\nref{prop:T8} with $(\mathbf X,\mathbf Y)=(\bfB,\mathbf 1)$, $\mathcal B={}_{\bfB}\mathcal C$, $F_{\mathbf 1}(W)=\bigoplus_{\mathbf m,\mathbf N}\underline{\mathrm{Hom}}(\mathbf N,\underline{\mathrm{End}}(\mathbf m)\otimes\mathbf N)=\bigoplus_{\mathbf m,\mathbf N}\underline{\mathrm{End}}(\mathbf m\otimes_{\bfB}\mathbf N)$; the labels match $f_{\mathbf m\mathbf N}$ by the argument of Proposition~\nref{prop:labels}. $\mathrm{Hom}(1,\underline{\mathrm{End}}(X))=\mathrm{End}(X)$. For $\mathbf m=\bfB$, $\bfB\otimes_{\bfB}\mathbf N=\mathbf N$. The last statement is the last part of Proposition~\ref{prop:cap}, with $\underline{\mathrm{End}}(X)\otimes i$ summands giving $\bar i\otimes\underline{\mathrm{End}}(X)\otimes i=\underline{\mathrm{End}}(X\otimes i)$.
\end{proof}
\emph{Status.} Proved here from the inputs of Theorem~\nref{thm:main}, Theorem~\ref{thm:through}, Lemma~\ref{lem:cardy} (which uses [KLM] and the correspondence between homomorphisms of extensions and $Q$-system morphisms), and one definition: the boundary-touching algebra of a configuration is generated by the boundary algebra and the unfolded defect, with strips glued through their phases. (H) is needed only to call $\underline{\mathrm{End}}(X)$ ``the boundary condition $X$''. The algebraic content is covered by the numerical checks (the $41$ triples with $K_{\bfC}=1$). The multiplicities $K_j\ge2$ are a prediction a naive fusion rule misses: the fused boundary carries the zero-mode algebra $\bigoplus M_{K_j}$ (Chan--Paton factors). The sum over $i$ in the last statement is the familiar statement that a Cardy strip carrying a Verlinde line $i$ shifts the boundary condition by $i$.

\section{Intertwiners between defects}\label{sec:sectors}
Point-like objects sitting on a defect line, or joining two different defect lines, are sectors in the sense of [BDH]. On the standard bicolored circle (white left arc, black right arc, colour points $p$ on top and $q$ at the bottom), a $D$--$E$-sector is a Hilbert space $H$ carrying representations of $\cA$ on intervals of the two arcs, of $D(I)$ for $I\ni p$ and of $E(I)$ for $I\ni q$. These are required to be compatible under restriction, with the top and bottom actions commuting on intervals with disjoint interiors. Vertical fusion of a $D$--$E$- and an $E$--$F$-sector is Connes fusion over $E$ of a half circle. Picture the circle around a point where an $E$-line coming from below turns into a $D$-line going up: the sector is the space of such junctions, including those that emit chiral charge into the bulk.

Besides L1--L4 we quote:
\begin{itemize}[leftmargin=2.6em]
\item[(L5)] [BDH] Gluing: compatible representations of the algebras of an open cover of a circle by intervals define a representation of the net on the circle.
\item[(L6)] Connes fusion over $\cA$ of a half circle of two DHR representations is the composition of the corresponding localized endomorphisms.
\end{itemize}

\begin{theorem}[sectors are bimodules]\label{thm:sectors}
Let $P,P'$ be $Q$-systems in $\mathcal Z$ localized in the strip. Then there is an equivalence of categories $\Sigma\colon\mathrm{Sect}(D_P,D_{P'})\to{}_P\mathcal Z_{P'}$ with $\Sigma(H_0(D_P))=P$, and $\Sigma(H\fus K)\cong\Sigma(H)\otimes_{P'}\Sigma(K)$ naturally in $H\in\mathrm{Sect}(D_P,D_{P'})$, $K\in\mathrm{Sect}(D_{P'},D_{P''})$.
\end{theorem}
\begin{proof}
\emph{From sectors to bimodules.} Insert the strip at both colour points; this gives a circle $S^\sharp$. The algebras $\cA(K)$, for $K\subset S^\sharp$ an interval inside an arc or inside some $J^\sharp$, cover $S^\sharp$. The sector axioms and naturality of $D_P$ under the structure maps make the restricted actions compatible, so by (L5) they glue to a representation $\pi$ of $\cA$ on $S^\sharp$. By complete rationality, $\pi$ is a direct sum of DHR representations; transport it to an object $X_H$ localized in an arc. The operator $m_\ell:=\rho^t(v_P)$ satisfies $m_\ell\pi(n)=\pi(\theta(n))m_\ell$, $m_\ell m_\ell=x\,m_\ell$ and $w^*m_\ell=1$, which are the axioms of a left $P$-module on $X_H$. Likewise $\rho^b(v_{P'})$ gives a right $P'$-module structure. The two commute because the top and bottom actions commute. Intertwiners of sectors are exactly the bimodule maps.

\emph{From bimodules to sectors.} Given $(\rho,m_\ell,m_r)\in{}_P\mathcal Z_{P'}$ with $\rho$ localized in an arc, let $\cA$ act through $\rho$, and send $v_P\mapsto$ the braiding-transport of $m_\ell$ to the top strip (as in $\alpha$-induction). The relations of $D_P(J)$ follow from the module axioms and the hexagon identities. For the bottom action, use convention C2: $D_{P'}(-I)=D_{P'}(I)^{\rm op}$, and a right $P'$-module is a left module over the opposite $Q$-system. So the bottom relations are the top computation applied to $P'^{\rm op}$ on the reflected circle. Top and bottom actions commute by the bimodule axiom, and functoriality is clear. These constructions are mutually inverse up to isomorphism, $H_0(D_P)=L^2(D_P(I))$ goes to $P$ with its regular bimodule structure, and $\Sigma$ is an equivalence.

\emph{Monoidality.} Let $M'=D_{P'}(J)=N'v'$ with $N'=\cA(J^\sharp)$, $J$ the half circle around $q$. \emph{(i) A one-element quasi-basis.} With $E'(nv')=d^{-1}nw'$, the element $b=d^{1/2}v'^*$ satisfies $bE'(b^*m)=m$ for all $m\in M'$. Indeed, for $m=nv'$,
\[
d\,v'^*E'(v'nv')=d\,v'^*E'(\theta'(n)x'v')=v'^*\theta'(n)x'w'=n\,v'^*x'w'.
\]
Using $v'^*=w'^*x'^*v'$, $v'a=\theta'(a)v'$, $x'^*\theta'(x')=x'x'^*$ (adjoint of the Frobenius identity $\theta'(x'^*)x'=x'x'^*$), $w'^*x'=1$ and $x'^*\theta'(w')=(\theta'(w'^*)x')^*=1$:
\[
v'^*x'w'=w'^*x'^*\theta'(x')\theta'(w')v'=w'^*x'x'^*\theta'(w')v'=v'.
\]
So $bE'(b^*m)=nv'=m$. \emph{(ii)} By (L4), $H\fus_{M'}K$ is the range of $q=d\,\rho_H(v'^*)\fus\lambda_K(v')$ on $H\fus_{N'}K$. \emph{(iii)} By (L6), $H\fus_{N'}K$ carries the object $X_H\otimes X_K$, with the outer actions of $D_P$, $D_{P''}$ on the outer factors. On it, $\lambda_K(v')$ is $l_K^*$ and $\rho_H(v'^*)$ is $r_H$, so $q=d\,(r_H\otimes1)(1\otimes l_K^*)$. Every module is equivalent to a standard one [BKLR, Lemma~3.22], and for a standard module $l_K^*$ is a module map $K\to P'\otimes K$ [BKLR, \S3.4]. So is $t_K:=(1\otimes l_K)(x'\eta'\otimes1)$, by the Frobenius identity. Under the free--forgetful adjunction $\mathrm{Hom}_{P'}(K,P'\otimes K)\cong\mathrm{Hom}(K,K)$, $f\mapsto(\varepsilon\otimes1)f$, both go to $1_K$: $(\varepsilon\otimes1)t_K=l_K(\eta\otimes1)=1$ and $(\varepsilon\otimes1)l_K^*=(l_K(\eta\otimes1))^*=1$. Hence $l_K^*=t_K$. (For the regular module this is $x'=(1\otimes m)(x'\eta'\otimes1)$, which fixes the normalization.) Therefore $q$ is a positive multiple of $(r_H\otimes l_K)(1\otimes x'\eta'\otimes1)$, i.e.\ of the action of the Casimir, whose normalized version $d^{-1}(r_H\otimes l_K)(1\otimes x'\eta'\otimes1)$ is an idempotent (cf.\ Lemma~\nref{lem:casimir}). A positive multiple of an idempotent that is itself a projection equals that idempotent, so $q$ is the normalized Casimir action. Its range is $X_H\otimes_{P'}X_K$ by the definition of the relative tensor product over a special Frobenius algebra [BKLR, \S4]. Naturality in $H$ and $K$ is clear from the construction.
\end{proof}
\emph{Status.} This closes the step that was left open in the earlier sketch of this result: identifying the Jones projection with the module idempotent. That step is now the explicit quasi-basis computation (i) together with the Casimir. Not written out: the compatibility of the monoidal structure with the associators, and the finiteness of sectors, which is assumed as in [BDH V].

\section{Junctions}\label{sec:junctions}
A junction is a point at which the configuration of defects and boundaries changes: two defects merge into one (a \emph{Y-junction}), a defect ends on a hard boundary at a point where the boundary condition changes (a \emph{T-junction}), or two defects cross (an \emph{X-junction}). Because transparent defects are topological for $\cA$, a junction can be moved and the lines tilted without changing anything, so the junction data reduce to the following: an intertwiner between the fused configuration on one side of the junction and the configuration on the other. \emph{Topological} junction fields are those commuting with all chiral observables, i.e.\ elements of relative commutants (zero-mode algebras); fields carrying a chiral charge $i$ are the corresponding charged intertwiners. Throughout, $f_{\mathbf m_1\mathbf m_2}$ and $D_{(\mathbf m_1,\mathbf m_2)}$ are as in Theorem~\nref{thm:main}, and for a projection $p$ in the relative commutant of a $Q$-system defect $D_P$, $pD_P$ denotes the corner $I\mapsto pD_P(I)p$ (with $\cA$ on monochrome intervals); it is the defect of the reduced $Q$-system of $p$ [Longo], [BKLR, \S4].

\begin{proposition}[Y-junctions]\label{prop:Y}
Let $X=\mathbf m_1\otimes_{\bfB}\mathbf m_2$ and $R=\mathrm{End}_{\bfA\text{-}\bfC}(X)$ the relative commutant of $D_{(\mathbf m_1,\mathbf m_2)}$ (Theorem~\nref{thm:main}(ii)). For every projection $p\in R$ onto a simple summand $\mathbf m\subset X$,
\[
pD_{(\mathbf m_1,\mathbf m_2)}\cong D_{\mathbf m}.
\]
Consequently the topological Y-junctions at which $\mathbf m_1$ and $\mathbf m_2$ merge into $\mathbf m$ are the elements of $\mathrm{Hom}_{\bfA\text{-}\bfC}(X,\mathbf m)\cong pRf_{\mathbf m_1\mathbf m_2}$, a space of dimension $N^{\mathbf m}_{\mathbf m_1\mathbf m_2}$, realized as operators in the relative commutant localized at the junction; their adjoints are the splitting junctions, and composing a splitting with a merging gives the matrix units of the block $M_{N^{\mathbf m}_{\mathbf m_1\mathbf m_2}}$ of $R$.
\end{proposition}
\begin{proof}
By Theorem~\nref{thm:kfold}(a) and Proposition~\nref{prop:labels}, the summand of $\Theta=Z[\bfA]\times Z[\bfB]\times Z[\bfC]$ cut out by $f_{\mathbf m_1\mathbf m_2}$ is $I_{\bfC}(\underline{\mathrm{End}}(X))$, and $R=\mathrm{Hom}(1,\underline{\mathrm{End}}(X))=\mathrm{End}(X)$. A projection $p\in\mathrm{End}(X)$ acts on $\underline{\mathrm{End}}(X)$ by composition on both sides, with $p\,\underline{\mathrm{End}}(X)\,p=\underline{\mathrm{End}}(pX)=\underline{\mathrm{End}}(\mathbf m)$ (functoriality of internal Hom), and $I_{\bfC}$, being additive and compatible with the lax structure, takes this corner to the corner of $I_{\bfC}(\underline{\mathrm{End}}(X))$ by $p$. Hence the reduced $Q$-system of $p$ is $I_{\bfC}(\underline{\mathrm{End}}(\mathbf m))$, which by Theorem~\nref{thm:kfold}(a) with $k=1$ is the summand $e_{\mathbf m}(Z[\bfA]\times Z[\bfC])$, the $Q$-system of $D_{\mathbf m}$. The two are isomorphic as algebras and $\underline{\mathrm{End}}(\mathbf m)$ is haploid, so Lemma~\nref{lem:Q5} upgrades this to an isomorphism of defects. The statement on junctions is the decomposition $\mathrm{End}(X)=\bigoplus_{\mathbf m}\mathrm{Hom}(\mathbf m,X)\otimes\mathrm{Hom}(X,\mathbf m)$.
\end{proof}
\emph{Numerical check} (script \texttt{junction\_check.py}). For $K_{\bfA}=K_{\bfB}=K_{\bfC}=A_4$ the pair $(3,3)$ has $N^{3}_{3,3}=2$. Taking a rank-one (non-central) idempotent $p$ in the $M_2$ block of the relative commutant, the corner $p\Theta p$ has dimension $108$ and the same graded character as the two-fold summand $e_3(L\times L)$, and differs from the other three two-fold summands, as Proposition~\ref{prop:Y} predicts.

\begin{proposition}[T-junctions and boundary-condition changes]\label{prop:T}
Let $\mathbf m\colon\bfA\to\bfB$, let $\mathbf N$ be a hard boundary condition of $\bfB$ and $\mathbf M$ one of $\bfA$. Consider a hard boundary of the phase $\bfA$ whose condition is $\mathbf M$ on one side of a point $p$, while on the other side the defect $\mathbf m$ emanates from $p$ into the bulk, with phase $\bfB$ and boundary condition $\mathbf N$ between the defect and the boundary. Under the dictionary of Theorem~\ref{thm:bcft}, the topological junction fields at $p$ form the space $\mathrm{Hom}_{\bfA}(\mathbf M,\mathbf m\otimes_{\bfB}\mathbf N)$ (or its dual, for the opposite time orientation), whose dimension is the multiplicity of $\mathbf M$ in $\mathbf m\otimes_{\bfB}\mathbf N$. Junction fields of chiral charge $i$ form $\mathrm{Hom}_{\bfA}(\mathbf M\otimes i,\mathbf m\otimes_{\bfB}\mathbf N)$.
\end{proposition}
\begin{proof}
Deform the defect towards the boundary on its side of $p$. By Theorem~\ref{thm:bcft} the boundary on that side becomes the boundary condition $Y':=\mathbf m\otimes_{\bfB}\mathbf N$, with chiral algebra $\underline{\mathrm{End}}(Y')$. The junction is then a boundary-condition-changing point between $\mathbf M$ and $Y'$. Both conditions live in the one chiral algebra $\underline{\mathrm{End}}(\mathbf M\oplus Y')$, whose charge-$i$ intertwiners are $\mathrm{Hom}(i,\underline{\mathrm{End}}(\mathbf M\oplus Y'))=\mathrm{Hom}_{\bfA}((\mathbf M\oplus Y')\otimes i,\mathbf M\oplus Y')$ (internal-Hom adjunction). The block from $\mathbf M$ to $Y'$ is the stated space. For $i=1$ these are elements of the zero-mode algebra $\mathrm{End}(\mathbf M\oplus Y')$, i.e.\ topological.
\end{proof}
For $\mathbf m=\bfB=\bfA$ this recovers the boundary-condition-changing fields of a single phase: topological ones exist only between isomorphic boundary conditions, and those of chiral charge $i$ form $\mathrm{Hom}_{\bfB}(\mathbf N\otimes i,\mathbf N')$, the usual Cardy-type multiplicities. This is consistent with the last statement of Proposition~\ref{prop:cap}, where the pinned Cardy strip produces $\bigoplus_i X\otimes i$. \emph{Status.} The algebraic statement is proved. Two things are not: that the deformation of the defect towards the boundary, applied on one side of $p$ only, is implemented by a unitary localized away from $p$ (which is what makes ``deform on one side'' legitimate), and the identification of the boundary-condition-changing fields of an LR boundary net with the charged intertwiners of its chiral algebra, which is the standard description and is quoted.

\begin{proposition}[X-junctions]\label{prop:X}
Let $\mathbf m_1\colon\bfA\to\bfB$ and $\mathbf m_2\colon\bfB\to\bfC$ cross at a point and continue as $\mathbf m_2'\colon\bfA\to\bfB'$, $\mathbf m_1'\colon\bfB'\to\bfC$. Put $X=\mathbf m_1\otimes_{\bfB}\mathbf m_2$ and $Y=\mathbf m_2'\otimes_{\bfB'}\mathbf m_1'$. Then every pair consisting of a merging junction $\mathbf m_1\otimes\mathbf m_2\to\mathbf m$ and a splitting junction $\mathbf m\to\mathbf m_2'\otimes\mathbf m_1'$ composes to a topological crossing. The crossings so obtained span
\[
\bigoplus_{\mathbf m}\mathrm{Hom}(X,\mathbf m)\otimes\mathrm{Hom}(\mathbf m,Y)=\mathrm{Hom}_{\bfA\text{-}\bfC}(X,Y),
\]
of dimension $\sum_{\mathbf m}N^{\mathbf m}_{\mathbf m_1\mathbf m_2}N^{\mathbf m}_{\mathbf m_2'\mathbf m_1'}$. Conversely, every junction between the two configurations that emits no transparent defect into the phase $\bfA$ is of this form.
\end{proposition}
\begin{proof}
On either side of the crossing the configuration fuses through its middle phase to $D_{(\mathbf m_1,\mathbf m_2)}$, resp.\ $D_{(\mathbf m_2',\mathbf m_1')}$. By Proposition~\ref{prop:Y} a merging vector selects a minimal projection $p$ of the block $\mathbf m$, with $pD_{(\mathbf m_1,\mathbf m_2)}\cong D_{\mathbf m}$, and a splitting vector selects $p'$, with $D_{\mathbf m}\cong p'D_{(\mathbf m_2',\mathbf m_1')}$. The composite is an isomorphism of corner defects, and isomorphisms compose, which gives the first statement. For the converse, the two defects are $D_{P_1}$, $D_{P_2}$ with $P_1=I_{\bfC}(\underline{\mathrm{End}}(X))$ and $P_2=I_{\bfC}(\underline{\mathrm{End}}(Y))$ (Theorem~\nref{thm:kfold}). By Theorem~\ref{thm:sectors} and the standard equivalence between bimodules and module functors [EGNO, \S7.11],
\[
\mathrm{Sect}(D_{P_1},D_{P_2})\simeq{}_{P_1}\mathcal Z_{P_2}\simeq\mathrm{Fun}_{\mathcal Z}\big(\mathrm{Mod}_{P_1}(\mathcal Z),\mathrm{Mod}_{P_2}(\mathcal Z)\big).
\]
By Proposition~\nref{prop:Q4} and the dominance remark after it, $\mathrm{Mod}_{P_i}(\mathcal Z)\simeq\mathrm{Mod}_{\underline{\mathrm{End}}(X_i)}(\mathcal D_{\bfC})$ as $\mathcal D_{\bfC}$-module categories, with $X_1=X$, $X_2=Y$. By Morita theory (proof of Theorem~\nref{thm:kfold}(c)) these are ${}_{\bfA}\mathcal C_{\bfC}$, via $M\mapsto M\otimes X_i$. Hence
\[
\mathrm{Sect}(D_{P_1},D_{P_2})\simeq\mathrm{Fun}_{\mathcal D_{\bfC}}({}_{\bfA}\mathcal C_{\bfC},{}_{\bfA}\mathcal C_{\bfC})\simeq\mathcal D_{\bfA}={}_{\bfA}\mathcal C_{\bfA},
\]
the last step by invertibility of ${}_{\bfA}\mathcal C_{\bfC}$. A sector is thus labelled by an $\bfA$--$\bfA$-bimodule $\mathbf a$, the transparent defect of the phase $\bfA$ emitted from the junction into the bulk. Under the Morita identifications the functor $\mathbf a\otimes_{\bfA}(-)$ is compared with the identity of ${}_{\bfA}\mathcal C_{\bfC}$ through $X$ and $Y$, and the multiplicity of the sector of type $\mathbf a$ in the junction is $\mathrm{Hom}_{\bfA\text{-}\bfC}(\mathbf a\otimes_{\bfA}X,Y)$ (Yoneda in the functor category). The junctions that emit nothing are those of type $\mathbf a=\bfA$, with junction space $\mathrm{Hom}_{\bfA\text{-}\bfC}(X,Y)$.
\end{proof}
\emph{Status.} The first statement (crossings built from Y-junctions, and their composition) is proved. The converse is proved from Theorem~\ref{thm:sectors}, Proposition~\nref{prop:Q4} and Morita theory. In addition it shows that the general junction between the two configurations also emits a transparent defect $\mathbf a$ of the phase $\bfA$, with junction space $\mathrm{Hom}(\mathbf a\otimes_{\bfA}X,Y)$; the topological crossings are the case $\mathbf a=\bfA$. The count says, in particular, when a defect can pass through another while the middle phase changes from $\bfB$ to $\bfB'$: exactly when $X$ and $Y$ share a simple summand.

\section{Intersecting transparent boundaries}\label{sec:crossing}
[BKLR] remark that ``it is much less clear which mathematical structure should be expected to describe situations where two transparent boundaries intersect each other.'' Proposition~\ref{prop:X} treated a crossing by assuming that it is an intertwiner between the fused configurations before and after it. This subsection removes that assumption. Crossing operators are characterized by Einstein causality alone, and the outcome is the following structure. Phases are the objects, boundary conditions are the 1-morphisms (composed by $\otimes_{\bfB}$), and crossings are 2-morphisms between composites of 1-morphisms. This is the 2-category of $Q$-systems in $\mathcal C$, bimodules and bimodule maps.

\textbf{Geometry.} Let $\ell,\ell'$ be timelike lines in two-dimensional Minkowski space crossing at the origin. They cut the plane into four open regions: $\Lambda_L,\Lambda_R$ to the left and right, and $\Lambda_-,\Lambda_+$ containing the negative and positive time axis. Since both lines are timelike, $\Lambda_\pm\subset V_\pm$, the open future and past light cones, while the wedges $W_L=\{x<-|t|\}$ and $W_R=\{x>|t|\}$ lie in $\Lambda_L$ and $\Lambda_R$. The causal complement of the origin is $W_L\cup W_R$. The phases are $\bfA$ on $\Lambda_L$, $\bfC$ on $\Lambda_R$, $\bfB$ on $\Lambda_-$ and $\bfB'$ on $\Lambda_+$. The boundary conditions are as in Proposition~\ref{prop:X}: $\mathbf m_1,\mathbf m_2$ on the past halves of the lines, $\mathbf m_2',\mathbf m_1'$ on the future halves. Put $X=\mathbf m_1\otimes_{\bfB}\mathbf m_2$ and $Y=\mathbf m_2'\otimes_{\bfB'}\mathbf m_1'$, and write $\mathcal M={}_{\bfA}\mathcal C_{\bfC}$.

\textbf{Past and future configurations.} Before the crossing the configuration is $\bfA\,|\,\bfB\,|\,\bfC$ with two parallel boundaries. Since the location of a transparent boundary is arbitrary [BKLR2, \S3.6], its fields are those of the universal construction $\mathcal D_-:=\cA_2\rtimes\Theta_f$, $\Theta_f:=f_{\mathbf m_1\mathbf m_2}(Z[\bfA]\times Z[\bfB]\times Z[\bfC])\cong I_{\bfC}(\underline{\mathrm{End}}(X))$ (Theorem~\nref{thm:kfold}, Proposition~\nref{prop:labels}). Likewise $\mathcal D_+:=\cA_2\rtimes\Theta'_{f'}$ after the crossing, with $\Theta'_{f'}\cong I_{\bfC}(\underline{\mathrm{End}}(Y))$. Both contain the \emph{outer system}, the fields of $\bfA$ and $\bfC$. It is the image of the extension by $Z[\bfA]\times Z[\bfC]$ under the algebra maps $\iota\colon Z[\bfA]\times Z[\bfC]\to\Theta_f$ and $\iota'\colon Z[\bfA]\times Z[\bfC]\to\Theta'_{f'}$ that insert the unit of the middle factor.

We quote:
\begin{itemize}[leftmargin=2.6em]
\item[(R)] [BKLR, Ch.~4], [LR95]. Consider representations of an extension $\cA_2\rtimes Q$, not necessarily local, whose restriction to $\cA_2$ is a sum of DHR representations. These correspond to left $Q$-modules in $\mathcal Z$, and intertwiners of $\cA_2$ and of the charged field correspond to module maps. The vacuum representation is the regular module $Q$, with $\mathrm{End}_Q(Q)=\mathrm{Hom}(1,Q)$ acting by right multiplication. Restricting a representation along an algebra map corresponds to restricting the module.
\end{itemize}
By Theorem~\nref{thm:kfold}(b), $\mathrm{Hom}(1,\Theta_f)=\mathrm{End}(X)$, so the vacuum representation of $\mathcal D_-$ splits as $\bigoplus_q\Theta_fq$ over a maximal family of orthogonal minimal projections $q\in\mathrm{End}(X)$. Two summands are equivalent iff their $q$'s lie in the same block. Write $H_-^{(\mu)}$ for the irreducible representation $\Theta_fq$ with $q$ in the block of a simple $\mu\in\mathcal M$, $N^\mu_X\neq0$, and similarly $H_+^{(\mu)}$ for $N^\mu_Y\ne0$. Seen from far away, the two boundaries look like a single $\bfA$--$\bfC$ boundary, and $\mu$ is its condition. I call $\mu$ the \emph{asymptotic channel}.

\begin{definition}[crossing field]\label{def:crossing}
A \emph{topological crossing field} from channel $\mu$ to channel $\mu'$ is a bounded operator $T\colon H_-^{(\mu)}\to H_+^{(\mu')}$ such that
\begin{enumerate}[nosep,label=(\roman*)]
\item $T\pi_-(a)=\pi_+(a)T$ for $a\in\cA_2(O)$ and all double cones $O$ (it is topological);
\item $T\pi_-(b)=\pi_+(b)T$ for $b\in(\cA_2\rtimes Z[\bfA])(O)$ with $O\subset W_L$, and for $b\in(\cA_2\rtimes Z[\bfC])(O)$ with $O\subset W_R$.
\end{enumerate}
\end{definition}
Condition (ii) is Einstein causality for an operator localized at the origin. It imposes nothing on the regions $\Lambda_\pm$, which are timelike to the origin, and these are exactly where the two descriptions differ.

\begin{lemma}[unit insertion]\label{lem:unitins}
Under the identifications of Theorem~\nref{thm:kfold}, $\iota=I_{\bfC}(\delta)$, where
\[
\delta\colon R:=\bigoplus_{\mathbf m\in\mathrm{Irr}\,\mathcal M}\underline{\mathrm{End}}(\mathbf m)\ \longrightarrow\ \underline{\mathrm{End}}(X)
\]
is the diagonal map. It is determined by the decomposition $X\cong\bigoplus_{\mathbf m}M^X_{\mathbf m}\otimes\mathbf m$, $M^X_{\mathbf m}=\mathrm{Hom}(\mathbf m,X)$: on $\mathrm{Hom}(W,-)$ it evaluates natural transformations $(-)\triangleleft W\Rightarrow\mathrm{id}_{\mathcal M}$ at $X$. On $\mathrm{Hom}(1,-)$ this gives $\iota(e_{\mathbf m})=z^X_{\mathbf m}$, the central projection of $\mathrm{End}(X)$ onto the $\mathbf m$-isotypic part.
\end{lemma}
\begin{proof}[Proof (object level)]
\emph{Step 1.} Let $\pi\colon I(1)\otimes Z\to I(F(Z))$ be the projection isomorphism, adjoint to $\varepsilon\otimes1$. Then $\pi\circ(\eta_{I(1)}\otimes1_Z)$ is adjoint to $(\varepsilon\circ F(\eta_{I(1)}))\otimes1=1_{F(Z)}$, because the unit of the algebra $I(1)$ is the adjunction unit at $1$. So $\pi\circ(\eta\otimes1)$ is the adjunction unit $u_Z\colon Z\to IF(Z)$.

\emph{Step 2.} Apply Step 1 with $Z=Z[\bfA]$ and $I=I_{\bfB}$: the map $Z[\bfA]\to Z[\bfB]\times Z[\bfA]\cong I_{\bfB}F_{\bfB}Z[\bfA]$, $x\mapsto\eta\otimes x$, is $u_{Z[\bfA]}$. Then $\Theta_f$ and $Z[\bfA]\times Z[\bfC]$ are identified with $I_{\bfC}F_{\bfC}(\cdot)$ by the projection isomorphism for $\bfC$, which is natural in its second variable and compatible with (T2). Under these identifications, $\iota=I_{\bfC}(F_{\bfC}(u_{Z[\bfA]}))$.

\emph{Step 3.} Use the description of (T4)--(T5) from the proof of Proposition~\nref{prop:T8}. The action of $I_{\bfB}F_{\bfB}Z[\bfA]$ is the free module functor on $F_{\bfB}Z[\bfA]\otimes(-)$, and $u$ acts as the unit of the free--forgetful adjunction. Transport this from ${}_{\bfB}\mathcal C_{\bfC}$ to $\mathcal M$ through the functors $\mathbf n'\otimes_{\bfB}(-)$, $\mathbf n'\colon\bfA\to\bfB$, which give $F_{\bfB}Z[\bfA]=\bigoplus_{\mathbf n'}\underline{\mathrm{End}}(\mathbf n')$ by Proposition~\nref{prop:T8}. The result is the evaluation of the action of $Z[\bfA]$ on $\mathcal M$ at the objects $\mathbf n'\otimes_{\bfB}\mathbf n$, which is $\delta$.

The $\mathrm{Hom}(1,-)$ statement follows because $\delta(1_{\mathbf m})$ is the identity on the $\mathbf m$-components of $X$.
\end{proof}
\emph{Status of the lemma.} Steps 1 and 2 are complete. Step 3 is the same Yoneda bookkeeping as Proposition~\nref{prop:T8}: it is described and not drawn, and its compatibility with products is part of Remark~\nref{rem:T8}. The consequence $\iota(e_{\mathbf m})f_{\mathbf m_1\mathbf m_2}=z^X_{\mathbf m}$ is a finite computation that the scripts of Appendix~\nref{app:num} could check; I have not run it. It is not a formality. An algebra map $\underline{\mathrm{End}}(\mathbf m'')\to\underline{\mathrm{End}}(\mu)$ between simple internal Ends can exist with $\mathbf m''\neq\mu$ (for instance the unit $1\to\underline{\mathrm{End}}(V)$ in $\mathrm{Rep}(S_3)$). So the lemma carries content: inside $D_\mu$, the outer fields are embedded in the standard way.

\begin{theorem}[crossings]\label{thm:crossing}
Assume the hypotheses of Theorem~\nref{thm:kfold} and Lemma~\ref{lem:unitins}. The space of topological crossing fields from channel $\mu$ to channel $\mu'$ is $0$ if $\mu\neq\mu'$. For $\mu=\mu'$ it is canonically
\[
\mathrm{Hom}_{\bfA\text{-}\bfC}(X,Y)=\bigoplus_{\mathbf m}\mathrm{Hom}(M^X_{\mathbf m},M^Y_{\mathbf m}),\qquad \dim=\sum_{\mathbf m}N^{\mathbf m}_{\mathbf m_1\mathbf m_2}N^{\mathbf m}_{\mathbf m_2'\mathbf m_1'}.
\]
The identification is compatible with composition: for a second crossing to a configuration with composite $Z$, composing crossing fields corresponds to composing in $\mathrm{Hom}(X,Y)\times\mathrm{Hom}(Y,Z)\to\mathrm{Hom}(X,Z)$.
\end{theorem}
\begin{proof}
\emph{Step 1: transport.} Let $O$ be any double cone and $O_0\subset W_L$. The charged field of $\bfA$ localized in $O$ is $v_O=u\,v_{O_0}$ with $u$ a charge transporter of $\cA_2$. By (i) and (ii), $Tv_O=uTv_{O_0}=uv_{O_0}'T=v'_OT$. So $T$ intertwines $\cA_2\rtimes Z[\bfA]$ on every double cone, and likewise for $\bfC$. The outer system is generated by these two subextensions (universal construction, [BKLR2, Prop.~5.1]). So $T$ intertwines the restrictions of $\pi_\pm$ along $\iota$, $\iota'$ to the outer system. Conversely, every such intertwiner satisfies (i) and (ii).

\emph{Step 2: modules.} By (R), the crossing fields are
\[
\mathrm{Hom}_{Z[\bfA]\times Z[\bfC]}\big(\iota^*(\Theta_fq),\iota'^*(\Theta'_{f'}q')\big).
\]
By Lemma~\ref{lem:unitins}, $\iota^*=I_{\bfC}(\delta)^*$. Both modules are $I_{\bfC}$ of modules in $\mathcal D_{\bfC}$, and by (T7) and Proposition~\nref{prop:Q4}, $I_{\bfC}$ is an equivalence $\mathrm{Mod}_R(\mathcal D_{\bfC})\simeq\mathrm{Mod}_{I_{\bfC}(R)}(\mathcal Z)$. So the space is
\[
\mathrm{Hom}_R\big(\delta^*(\underline{\mathrm{End}}(X)q),\ \delta'^*(\underline{\mathrm{End}}(Y)q')\big).
\]

\emph{Step 3: decomposition.} Right multiplication by $q$ gives $\underline{\mathrm{End}}(X)q=\underline{\mathrm{Hom}}(\mu,X)$, with $\underline{\mathrm{End}}(X)$ acting by postcomposition. Through the diagonal map, $\delta^*\underline{\mathrm{Hom}}(\mu,X)\cong\bigoplus_{\mathbf m}M^X_{\mathbf m}\otimes\underline{\mathrm{Hom}}(\mu,\mathbf m)$, with $R$ acting on the summand $\mathbf m$ through $\underline{\mathrm{End}}(\mathbf m)$. The same holds for $Y$ and $\mu'$.

\emph{Step 4: Ostrik.} $\mathcal M$ is indecomposable and every simple $\mathbf m$ generates it. So $x\mapsto\underline{\mathrm{Hom}}(x,\mathbf m)$ is an equivalence of $\mathcal M^{\rm op}$ with left $\underline{\mathrm{End}}(\mathbf m)$-modules in $\mathcal D_{\bfC}$ [Ostrik], [EGNO, Thm.~7.10.1]. Hence $\mathrm{Hom}_{\underline{\mathrm{End}}(\mathbf m)}(\underline{\mathrm{Hom}}(\mu,\mathbf m),\underline{\mathrm{Hom}}(\mu',\mathbf m))\cong\mathrm{Hom}_{\mathcal M}(\mu',\mu)=\delta_{\mu\mu'}\mathbb C$. Summing over $\mathbf m$ with the multiplicity spaces gives the formula, and Schur's lemma identifies the sum with $\mathrm{Hom}(X,Y)$. Every identification used is a functor, so composition is respected.
\end{proof}

\begin{corollary}[channels match across a crossing]\label{cor:channels}
Suppose a single theory realizes the crossing: one Hilbert space $H$ carrying representations $\pi_\pm$ of $\mathcal D_\pm$ of vacuum type that coincide on the outer fields localized in $W_L\cup W_R$. Then $\pi_-(z^X_{\mathbf m})=\pi_+(z^Y_{\mathbf m})$ for every $\mathbf m$. In particular a channel $\mathbf m$ with $N^{\mathbf m}_X\neq0$ but $N^{\mathbf m}_Y=0$ is absent from $H$. If $X$ and $Y$ share no simple summand, the two boundaries cannot cross without emitting a defect.
\end{corollary}
\begin{proof}
By Step~1 of the theorem, $\pi_-$ and $\pi_+$ coincide on the whole outer system. By Lemma~\ref{lem:unitins} the outer element $\iota(e_{\mathbf m})=\iota'(e_{\mathbf m})$ is $z^X_{\mathbf m}$ in the past description and $z^Y_{\mathbf m}$ in the future one.
\end{proof}

\begin{remark}[the structure, and how it relates to the earlier results]\label{rem:2cat}
(a) \emph{The answer to the question of [BKLR].} Intersecting transparent boundaries are described by the 2-category $\mathrm{Bim}(\mathcal C)$: simple $Q$-systems (phases), bimodules (boundary conditions) composed by $\otimes_{\bfB}$, and bimodule maps. Two boundaries do not have a canonical crossing. A crossing is a choice of vector in $\mathrm{Hom}_{\bfA\text{-}\bfC}(X,Y)$, and when this space is $0$ the boundaries cannot cross. In particular no braiding or other extra structure is involved. Theorem~\ref{thm:crossing} reproduces the count of Proposition~\ref{prop:X}, now derived from causality rather than assumed, and Proposition~\ref{prop:Y} decomposes every crossing into a merge followed by a split.

(b) \emph{Compositions.} Vertical composition, meaning crossings in sequence in time, is the composition statement of Theorem~\ref{thm:crossing}. Horizontal composition, meaning crossings side by side or a crossing next to a parallel boundary, and the interchange law are treated in Section~\ref{sec:horizontal} (Corollary~\ref{cor:interchange}).

(c) \emph{Emitted defects.} A transparent $\bfA$--$\bfA$ defect $\mathbf a$ leaving the crossing into $\Lambda_L$ must avoid $W_L$, which contains no timelike ray from the origin. The future configuration is then $\mathbf a\otimes_{\bfA}Y$, and the theorem gives $\mathrm{Hom}(X,\mathbf a\otimes_{\bfA}Y)$. This is the dual form of the count $\mathrm{Hom}(\mathbf a\otimes_{\bfA}X,Y)$ in Proposition~\ref{prop:X}.

(d) \emph{Context.} This is the structure used for defect junctions in the TFT approach to rational CFT [FRS]. What is new here is that it follows from locality of the operator-algebraic boundary nets.

(e) \emph{Not addressed.} We characterize the crossing fields between two given asymptotic configurations. We do not construct, from scratch, a net on regions containing the crossing that realizes a given crossing field, and we do not treat intersections of hard boundaries.
\end{remark}

\section{Horizontal composition: crossings at spacelike separation}\label{sec:horizontal}
Theorem~\ref{thm:crossing} gives the 2-morphisms and their vertical composition. In a 2-category, horizontal composition puts two 2-morphisms side by side, $\alpha\colon X\Rightarrow Y$ in $\mathrm{Bim}(\bfA,\bfC)$ and $\beta\colon X'\Rightarrow Y'$ in $\mathrm{Bim}(\bfC,\bfE)$ giving $\alpha\otimes_{\bfC}\beta$. Physically this means two crossings at spacelike-separated points with the phase $\bfC$ between them. The interchange law then says that it does not matter which crossing happens first. This is the kind of statement Einstein causality should give, and the following shows that it does. The route avoids the fusion of sectors, and with it the associator coherence left open after Theorem~\ref{thm:sectors}.

\textbf{Configurations with several crossings.} Let $p_1,\dots,p_r$ be pairwise spacelike points, each the crossing point of two transparent boundaries as in Section~\ref{sec:crossing}, and allow further boundaries that cross nothing (\emph{spectators}). A phase is \emph{visible} if its region meets the causal complement of $\{p_1,\dots,p_r\}$, and \emph{hidden} otherwise. The hidden phases are exactly the phases $\bfB$ (before) and $\bfB'$ (after) of the individual crossings, whose regions lie in $V_\mp(p_j)$. Order the visible phases from left to right as $\mathbf P_0,\dots,\mathbf P_k$; they are the same before and after. Between $\mathbf P_{i-1}$ and $\mathbf P_i$ the boundaries before the crossings compose to $X_i$, a single simple bimodule for a spectator and a composite such as $\mathbf m_1\otimes_{\bfB}\mathbf m_2$ at a crossing; after the crossings they compose to $Y_i$. Put $X_{\rm tot}=X_1\otimes\cdots\otimes X_k$ and $Y_{\rm tot}=Y_1\otimes\cdots\otimes Y_k$ (relative tensor products over the visible phases, bracketed as in Theorem~\nref{thm:kfold}), and $\mathcal M={}_{\mathbf P_0}\mathcal C_{\mathbf P_k}$. As before, the past and future fields are the universal constructions $\cA_2\rtimes\Theta_f$, $\cA_2\rtimes\Theta'_{f'}$ with $\Theta_f\cong I(\underline{\mathrm{End}}(X_{\rm tot}))$ and $\Theta'_{f'}\cong I(\underline{\mathrm{End}}(Y_{\rm tot}))$, where $I=I_{\mathbf P_k}$. Their irreducible vacuum-type representations $H_\mp^{(\mu)}$ are labelled by asymptotic channels $\mu\in\mathrm{Irr}\,\mathcal M$. The \emph{visible system} is the image of $Z[\mathbf P_0]\times\cdots\times Z[\mathbf P_k]$ under the maps $\iota,\iota'$ that insert the units of the hidden factors. A \emph{crossing field at $\{p_j\}$} is a bounded $T\colon H_-^{(\mu)}\to H_+^{(\mu')}$ that intertwines $\cA_2$ everywhere and intertwines the fields of each visible phase on double cones in the causal complement of $\{p_j\}$ (Definition~\ref{def:crossing} with $W_L\cup W_R$ replaced by that causal complement).

\begin{lemma}[unit insertion, several factors]\label{lem:unitins2}
Under the identifications of Theorem~\nref{thm:kfold}, $\iota=I(\delta)$. Here $\delta\colon R_{\rm vis}:=\bigoplus_t\underline{\mathrm{End}}(N_t)\to\underline{\mathrm{End}}(X_{\rm tot})$ runs over tuples $t=(\mathbf n_1,\dots,\mathbf n_k)$ of simple $\mathbf P_{i-1}$--$\mathbf P_i$-bimodules, $N_t=\mathbf n_1\otimes\cdots\otimes\mathbf n_k$, and $\delta$ is the diagonal map determined by the decompositions $X_i\cong\bigoplus_{\mathbf n}M^{X_i}_{\mathbf n}\otimes\mathbf n$, with $X_{\rm tot}\cong\bigoplus_t(\bigotimes_iM^{X_i}_{\mathbf n_i})\otimes N_t$.
\end{lemma}
\begin{proof}
Insert the hidden units one at a time. For each insertion, Steps~1 and~2 of the proof of Lemma~\ref{lem:unitins} hold with $Z[\bfA]$ replaced by the product $Z$ of the factors to its left (they hold for any algebra $Z$). Step~3 uses only the description in Proposition~\nref{prop:T8} of $F I(V)$ for $V=F(Z)=\bigoplus\underline{\mathrm{End}}(X_{t'})$ (Theorem~\nref{thm:kfold}). Factors to the right of the insertion are added afterwards by $I_jF_j$, which is a functor and so carries the diagonal map along.
\end{proof}

\begin{theorem}[locality of crossings]\label{thm:horizontal}
Assume the hypotheses of Theorem~\nref{thm:kfold} and Lemma~\ref{lem:unitins2}.
\begin{enumerate}[nosep,label=(\alph*)]
\item The crossing fields at $\{p_j\}$ from channel $\mu$ to channel $\mu'$ form $0$ if $\mu\ne\mu'$. For $\mu=\mu'$ they are exactly the operators that, under the module picture $H_\mp^{(\mu)}\leftrightarrow I(\underline{\mathrm{Hom}}(\mu,X_{\rm tot}))$, resp.\ $I(\underline{\mathrm{Hom}}(\mu,Y_{\rm tot}))$, are postcomposition with an element of
\[
\mathrm{Hom}(X_1,Y_1)\otimes\cdots\otimes\mathrm{Hom}(X_k,Y_k)\ \xrightarrow{\ \otimes\ }\ \mathrm{Hom}_{\mathcal M}(X_{\rm tot},Y_{\rm tot}).
\]
This map is injective, so the crossing fields form $\bigotimes_i\mathrm{Hom}(X_i,Y_i)$. A spectator gap contributes $\mathrm{End}(\mathbf m)=\mathbb C$.
\item Composition of crossing fields is composition of the corresponding elements of $\mathrm{Hom}(X_{\rm tot},Y_{\rm tot})$.
\end{enumerate}
\end{theorem}
\begin{proof}
\emph{(i) Reduction to modules.} The transport argument of Theorem~\ref{thm:crossing} shows that $T$ intertwines the fields of each visible phase on all double cones. So, by (R), $T$ is a map of modules over the visible system. By Lemma~\ref{lem:unitins2} and Proposition~\nref{prop:Q4} it is a map of $R_{\rm vis}$-modules $\delta^*\underline{\mathrm{Hom}}(\mu,X_{\rm tot})\to\delta'^*\underline{\mathrm{Hom}}(\mu',Y_{\rm tot})$ in $\mathcal D_{\mathbf P_k}$, where $\underline{\mathrm{End}}(X_{\rm tot})q=\underline{\mathrm{Hom}}(\mu,X_{\rm tot})$ as in Theorem~\ref{thm:crossing}.

\emph{(ii) Counting.} As $R_{\rm vis}$-modules,
\[
\delta^*\underline{\mathrm{Hom}}(\mu,X_{\rm tot})\cong\bigoplus_t\Big(\bigotimes_iM^{X_i}_{\mathbf n_i}\Big)\otimes\underline{\mathrm{Hom}}(\mu,N_t),
\]
with $R_{\rm vis}$ acting on the summand $t$ through $\underline{\mathrm{End}}(N_t)$. Distinct $t$ give distinct algebra summands, so the Hom space splits over $t$. Each $N_t$ is non-zero (a relative tensor product of simple bimodules between simple $Q$-systems), so it generates the indecomposable $\mathcal M$. Hence Ostrik's theorem as in Theorem~\ref{thm:crossing} gives
\[
\mathrm{Hom}_{\underline{\mathrm{End}}(N_t)}\big(\underline{\mathrm{Hom}}(\mu,N_t),\underline{\mathrm{Hom}}(\mu',N_t)\big)\cong\mathrm{Hom}_{\mathcal M}(\mu',\mu).
\]
Summing over $t$, the space is $\delta_{\mu\mu'}\bigoplus_t\bigotimes_i\mathrm{Hom}(M^{X_i}_{\mathbf n_i},M^{Y_i}_{\mathbf n_i})=\delta_{\mu\mu'}\bigotimes_i\mathrm{Hom}(X_i,Y_i)$.

\emph{(iii) Postcomposition.} Let $\gamma=\gamma_1\otimes\cdots\otimes\gamma_k$ with $\gamma_i\in\mathrm{Hom}(X_i,Y_i)$. By Schur, $\gamma_i=\bigoplus_{\mathbf n}(\gamma_i)_{\mathbf n}\otimes1_{\mathbf n}$, so $\gamma$ preserves the summands $t$ and is the identity on each $N_t$. Hence postcomposition with $\gamma$ commutes with $\delta(R_{\rm vis})$ and is an $R_{\rm vis}$-module map. The functor $\underline{\mathrm{Hom}}(\mu,-)\colon\mathcal M\to\mathrm{Mod}\text{-}\underline{\mathrm{End}}(\mu)$ is an equivalence, hence faithful. So $\gamma\mapsto$ postcomposition is injective on the image of $\otimes$. It is also injective on the tensor product itself, because the components for different $t$ land in different blocks. By the count in (ii) it is onto.

\emph{(b)} Postcomposition is functorial.
\end{proof}
For $k=1$ with the two lines crossing (visible phases $\bfA,\bfC$), (a) is Theorem~\ref{thm:crossing}, and it also identifies the isomorphism there as postcomposition.

\begin{corollary}[horizontal composition and interchange]\label{cor:interchange}
Let $p,p'$ be spacelike, with $p$ to the left. At $p$, boundaries composing to $X$ cross into $Y$ (phases $\bfA\to\bfC$); at $p'$, boundaries composing to $X'$ cross into $Y'$ (phases $\bfC\to\bfE$). There are four configurations, labelled by which crossings have happened: $(--)$, $(+-)$, $(-+)$, $(++)$. For $\alpha\in\mathrm{Hom}(X,Y)$ and a spectator $S$ in the other gap, let $T^{S}_\alpha$ be the crossing field at $p$ given by $\alpha\otimes1_S$ (Theorem~\ref{thm:horizontal}(a)), and define $T^S_\beta$ at $p'$ similarly. Then:
\begin{enumerate}[nosep,label=(\alph*)]
\item \emph{Whiskering.} $T^S_\alpha$ is the unique crossing field at $p$ with label $\alpha$, for every spectator $S$. The label does not depend on what happens elsewhere.
\item \emph{Interchange.} As operators $H_{--}^{(\mu)}\to H_{++}^{(\mu)}$,
\[
T^{Y}_\beta\,T^{X'}_\alpha\;=\;T^{Y'}_\alpha\,T^{X}_\beta\;=\;T_{\alpha\otimes_{\bfC}\beta},
\]
the crossing field at $\{p,p'\}$ with label $\alpha\otimes_{\bfC}\beta\in\mathrm{Hom}(X,Y)\otimes\mathrm{Hom}(X',Y')$.
\item \emph{Horizontal composition.} The crossing fields at $\{p,p'\}$ are exactly the linear combinations of the $T_{\alpha\otimes\beta}$. The phase $\bfC$ between the crossings is visible, and this forces the label to factorize: crossing fields at the pair form the subspace $\mathrm{Hom}(X,Y)\otimes\mathrm{Hom}(X',Y')$ of $\mathrm{Hom}(X\otimes_{\bfC}X',Y\otimes_{\bfC}Y')$. The whole of that Hom space is the crossing space of the configuration fused across $\bfC$ (Theorem~\ref{thm:crossing} with $\bfC$ hidden).
\end{enumerate}
\end{corollary}
\begin{proof}
(a) For a crossing at $p$ in the configuration $(--)\to(+-)$, the visible gaps are $X\to Y$ at $p$ and the spectator boundaries of the configuration at $p'$, each contributing $\mathbb C$. So Theorem~\ref{thm:horizontal}(a) gives $\mathrm{Hom}(X,Y)$, realized as postcomposition with $\alpha\otimes1_S$.

(b) Both composites are crossing fields at $\{p,p'\}$: each factor intertwines the fields of $\bfA$, $\bfC$, $\bfE$ in the causal complement of its own point, which contains the causal complement of $\{p,p'\}$. By Theorem~\ref{thm:horizontal}(b) they are postcomposition with $(1_Y\otimes\beta)(\alpha\otimes1_{X'})$, resp.\ $(\alpha\otimes1_{Y'})(1_X\otimes\beta)$. Both equal $\alpha\otimes\beta$ by bifunctoriality of $\otimes_{\bfC}$.

(c) This is Theorem~\ref{thm:horizontal}(a) with visible phases $\bfA,\bfC,\bfE$.
\end{proof}

\begin{remark}[coherence, and what this achieves]\label{rem:horizontal}
(a) \emph{Associators.} Each configuration is identified with a single object $X_{\rm tot}$, using one fixed bracketing (that of Theorem~\nref{thm:kfold}). A different bracketing changes the identification by the associator of $\mathrm{Bim}(\mathcal C)$, under which $\gamma$ becomes its conjugate by that associator. The operators themselves are defined physically and do not depend on any choice. So the labels of crossing fields compose with the coherence of $\mathrm{Bim}(\mathcal C)$, which holds by Mac Lane, and no separate coherence argument is needed. The input is the associativity of fusion and of the braided product already used in Theorem~\nref{thm:kfold} and Proposition~\nref{prop:assoc}. In particular, Theorem~\ref{thm:sectors} and its unwritten associator coherence are not used here.

(b) \emph{The 2-category is realized by operators.} Phases, boundary conditions and crossing fields, with vertical composition as composition in time (Theorem~\ref{thm:crossing}) and horizontal composition as spacelike juxtaposition (Corollary~\ref{cor:interchange}), form $\mathrm{Bim}(\mathcal C)$. Units are the trivial crossings, and the interchange law is an operator identity that follows from Einstein causality. This completes the answer of Remark~\ref{rem:2cat} to the question of [BKLR], within the hypotheses stated.

(c) \emph{Not covered.} Crossings at lightlike separation; crossings whose hidden regions are not contained in the light cones of the crossing points (e.g.\ boundaries tilted to null, Section~\ref{sec:null}); and, as before, a construction of the crossing net from scratch.
\end{remark}

\section{Corners}\label{sec:corners}
In Lorentzian signature a hard boundary is a timelike curve, and a \emph{corner} is a point where it bends non-smoothly, so that its velocity jumps. A corner may coincide with a change of boundary condition or with the end of a transparent defect. In light-cone coordinates the bulk is $\{u>\varphi(v)\}$ for an increasing function $\varphi$. The straight boundary is $\varphi=\mathrm{id}$, and a corner at $v_0$ is a jump of $\varphi'$ at $v_0$, with ratio $\lambda=\varphi'(v_0^+)/\varphi'(v_0^-)\ne1$.

\textbf{Smooth bends are invisible.} The bulk nets are covariant under independent reparametrizations $(u,v)\mapsto(f(u),g(v))$, $f,g\in\mathrm{Diff}_+$, since both chiral factors of $\cA_2$ are diffeomorphism covariant and so are full-centre extensions. For smooth $\varphi$ the map $(u,v)\mapsto(u,\varphi^{-1}(v))$ takes the boundary $u=\varphi(v)$ to $u=v$. It carries the LR boundary net of the straight boundary to a boundary net for the curved one, with the same chiral algebra $Y$ and the same bulk, and all statements of this section transfer verbatim. So only non-smooth points carry information.

\textbf{Corners as chiral defects.} Under the unfolding of Lemma~\ref{lem:unfold}, a corner at $v_0$ becomes a point $k$ of the light ray. On either side of $k$ the reparametrization is smooth, but the two smooth pieces $\psi_-$ (identity) and $\psi_+$ meet at $k$ with derivative ratio $\lambda$. Define the \emph{kink defect} of the chiral net $A$ by
\[
K_\lambda(I):=A(I_\circ)\vee A(\psi_+(I_\bullet))=A\big(I_\circ\cup\psi_+(I_\bullet)\big)\qquad(I\ni k),
\]
with $A(J)$ on white intervals and structure maps $\alpha_{\psi_+}$ on black intervals. The equality is strong additivity, using $\psi_+(k)=k$.


\begin{lemma}[transport along a kinked homeomorphism]\label{lem:transport}
Let $S$ be a bicolored circle with colour points $p,q$, and let $\Psi$ be an orientation-preserving homeomorphism of $S$ fixing $p$ and $q$, smooth on $S\setminus\{p,q\}$. In the local coordinates at $p$ and at $q$, assume $\Psi$ is the identity on the white side and the dilation $x\mapsto\lambda_px$, resp.\ $\lambda_qx$, on the black side near the colour point. Define $(\Psi^*\Id_A)(I):=A(\Psi(I))$ and, for a morphism $f$, $(\Psi^*\Id_A)(f):=A(\Psi f\Psi^{-1})$.
\begin{enumerate}[nosep,label=(\roman*)]
\item $\Psi f\Psi^{-1}$ is smooth for every morphism $f$ of bicolored intervals (orientation-preserving or reversing, colour- and coordinate-preserving), so $\Psi^*\Id_A$ is a well-defined functor, and it satisfies all the defect axioms of Definition~\nref{def:defect}, including the vacuum-sector axiom (4).
\item For a $D$--$E$-sector $(H,\rho)$ of defects of this type, $(H,\rho\circ\Psi^*)$ is a $\Psi^*D$--$\Psi^*E$-sector. If $\Psi$ is smooth on the half circle over which two sectors are fused, then $\Psi^*(H\fus K)\cong\Psi^*H\fus\Psi^*K$.
\end{enumerate}
\end{lemma}
\begin{proof}
(i) Morphisms of bicolored intervals preserve the local coordinate, so they are the identity near the colour point. Hence $\Psi f\Psi^{-1}$ is the identity near the colour point, and it is smooth elsewhere because $\Psi$ and $f$ are. For an orientation-reversing $f$ exchanging $p$ and $q$ in the coordinates, $\Psi f\Psi^{-1}$ is the identity in the coordinates when $\lambda_p=\lambda_q$, which is the case needed below. Isotony, locality and strong additivity are those of $A$, because $\Psi$ maps intervals to intervals and disjoint intervals to disjoint ones. For (4): $L^2((\Psi^*\Id_A)(I))=L^2(A(\Psi I))$; the left action is that of $A(\Psi I)$, and the right action is that of $A(\Psi I')=A(\hat\jmath\,\Psi I)$ through $A(\hat\jmath)$, with $\hat\jmath:=\Psi j\Psi^{-1}$ a smooth orientation-reversing involution fixing $\partial(\Psi I)$. For a monochrome $J\subset I$ touching $\partial I$ and $\bar J=j(J)$ we have $\Psi(\bar J)=\hat\jmath(\Psi J)$. The required extension to $A(\Psi(J\cup\bar J))=A(\Psi J\cup\hat\jmath\Psi J)$ is the vacuum-sector property of the net $A$ for the interval $\Psi I$ and the involution $\hat\jmath$, which holds for every smooth involution because $H_0(\hat\jmath)$ is defined for all of $\mathrm{Conf}_-$ (Section~\nref{sec:nets}).
(ii) The compatibility and commutation conditions are preserved, because $\Psi$ preserves inclusions and disjointness. If $\Psi$ is smooth on the half circle $S_\perp$ (and preserves it), then $\alpha_\Psi$ restricted to $A(S_\perp)$ is implemented by $U(\phi)$ for any $\phi\in\mathrm{Diff}_+$ agreeing with $\Psi$ on $S_\perp$. Connes fusion is natural under isomorphisms of the middle algebra that act on both modules.
\end{proof}

\begin{proposition}[corners]\label{prop:corner}
\begin{enumerate}[nosep,label=(\alph*)]
\item $K_\lambda$ is an $A$--$A$ defect, and the configuration ``hard boundary with chiral algebra $Y$ and a corner of ratio $\lambda$ at the same point'' unfolds to $K_\lambda\cfus D^{\rm ch}_Y\cong D^{\rm ch}_{Y}\cfus K_\lambda$. Corners factorize from boundary conditions, and in particular from all the topological data of the preceding subsections (boundary conditions, T-junction spaces).
\item For every conformal net $A$ and every $\lambda\ne1$, $K_\lambda\not\cong\Id_A$ as defects: the kinked homeomorphism cannot be implemented on the local algebras of intervals containing the corner.
\item $K_\lambda$ is equivalent to $\Id_A$ through an invertible sector, whose arc representation is the irreducible, index-one, proper soliton $\sigma_\lambda$ of [DIT]. Solitons with different $\lambda$ are inequivalent [DIT]. Hence a corner is invisible at the level of equivalence classes of defects, and it carries a charge: the class of the soliton $\sigma_\lambda$, labelled by the velocity jump $\lambda\in\mathbb R_+$.
\end{enumerate}
\end{proposition}
\begin{proof}
(a) Isotony, locality and strong additivity hold because $K_\lambda(I)$ is $A$ of an interval and the structure maps are restrictions of covariance automorphisms. On the circle, $K_\lambda=\chi^*\Id_A$ for a homeomorphism $\chi$ of the type of Lemma~\ref{lem:transport} with $\lambda_p=\lambda_q=\lambda$, so all the defect axioms, including the vacuum-sector axiom, hold by Lemma~\ref{lem:transport}(i). For the factorization: $Y$ is localized in a strip away from $k$, and transporting it across $k$ replaces $\theta_Y$ by $\alpha_{\psi_+}\theta_Y\alpha_{\psi_+}^{-1}$. By covariance of DHR sectors this is a localized endomorphism unitarily equivalent to $\theta_Y$, with the covariance cocycle as the unitary, and it carries an equivalent $Q$-system. Equivalent $Q$-systems have isomorphic defects. The topological data were computed from $Y$ alone.

(b) Choose $\psi_+$ equal to the identity near $\infty$, so that $\chi_\lambda$ (identity left of $k$, $\psi_+$ right of $k$) is a homeomorphism of the circle, smooth except at $k$, where its one-sided derivatives have ratio $\lambda$. This is the class $\mathrm{Diff}_+(S^1,k)$ of [DIT, \S3.2], rotated so that the singular point is $k$. [DIT] define the soliton $\sigma_{\chi}$ on intervals $J\not\ni k$ by $\sigma_\chi|_{A(J)}=\mathrm{Ad}\,U(\chi_J)$, with $\chi_J\in\mathrm{Diff}_+(S^1)$ agreeing with $\chi$ on $J$. They prove (Thm.~3.6) that for $r(\chi)=\lambda\ne1$ it is proper, i.e.\ does not extend to a representation of the net on the whole circle. Suppose $\beta\colon\Id_A\to K_\lambda$ were a natural isomorphism of defects. On intervals $I\ni k$ it gives automorphisms $\beta_I$ of $A(I)$. Naturality with respect to the structure maps says $\beta_I|_{A(J)}=\mathrm{id}=\sigma_\chi|_{A(J)}$ for white $J\subset I$, $\beta_I|_{A(J)}=\alpha_{\psi_+}=\sigma_\chi|_{A(J)}$ for black $J\subset I$, and $\beta_J=\beta_I|_{A(J)}$ for $k\in J\subset I$. Thus $\{\beta_I\}_{I\ni k}\cup\{\sigma_\chi|_{A(J)}\}_{J\not\ni k}$ is a compatible family on all intervals of the circle, i.e.\ an extension of $\sigma_\chi$ to the whole net. This contradicts properness. The argument behind [DIT, Thm.~3.6] is modular: an extension would make a one-sided dilation $\mathrm{Ad}\,U(\delta(t))$, $t=-\log\lambda$, inner in the half-line algebra, whose modular group it is by Bisognano--Wichmann; this is impossible for $t\ne0$ in a type III$_1$ factor.

\emph{Illustration (free fermion).} Here the obstruction is visible directly: the Bogoliubov automorphism of $\chi_\lambda$ would require the localized difference of the Hardy projection $P$ and $VPV^*$, $(Vf)(s)=\chi_\lambda'(s)^{1/2}f(\chi_\lambda^{-1}(s))$, to be Hilbert--Schmidt (Powers--St{\o}rmer--Araki). Near $k$, with $s=k-a$, $t=k+b$, the kernel difference is $\sqrt\lambda/(a+\lambda b)-1/(a+b)$, homogeneous of degree $-1$, so its Hilbert--Schmidt norm diverges logarithmically unless $\lambda=1$.

(c) Let $H$ be the vacuum Hilbert space of $A$ on the circle. Take the top action to be that of $K_\lambda$, the white arc and the bottom colour point untwisted, and the black arc carrying the soliton $\sigma$ of a diffeomorphism $\tilde\psi$ of the closed black arc with $\tilde\psi=\psi_+$ near $p$ and $\tilde\psi=\mathrm{id}$ near $q$. Such a diffeomorphism exists, since the endpoint derivative is free. This is a soliton of the type of [DIT, \S3.2] with ratio $\lambda$ at $p$, and the compatibilities hold by construction. Equivalently, $H=\Psi^*H_0(\Id_A)$ for the homeomorphism $\Psi$ that is the identity on the white arc and $\tilde\psi$ on the black arc. It has a kink only at $p$, with $\lambda_p=\lambda$ and $\lambda_q=1$. By Lemma~\ref{lem:transport}(ii) it is a $K_\lambda$--$\Id_A$-sector satisfying all the sector axioms of [BDH]. \emph{Invertibility.} Let $\bar H=\bar\Psi^*H_0(\Id_A)$ be the $\Id_A$--$K_\lambda$-sector with the kink at $q$. Vertical fusion is over the half circle containing $q$, resp.\ $p$, where the other sector's $\Psi$ is smooth, so Lemma~\ref{lem:transport}(ii) gives $H\fus\bar H\cong\Psi_2^*(H_0(\Id_A)\fus H_0(\Id_A))\cong\Psi_2^*H_0(\Id_A)=H_0(K_\lambda)$, with $\Psi_2$ kinked at both points. The same argument gives $\bar H\fus H\cong H_0(\Id_A)$. So $K_\lambda\simeq\Id_A$. Independently, [DIT, Prop.~3.5 and \S3.1] show that the arc representation is irreducible, of index $1$ and of positive energy. Its class is that of $\sigma_\lambda$, and [DIT, Thm.~3.6] gives $\sigma_{\lambda_1}\simeq\sigma_{\lambda_2}$ iff $\lambda_1=\lambda_2$.
\end{proof}

\emph{Interpretation.} (b) is the Lorentzian, operator-algebraic counterpart of the corner anomaly of Euclidean boundary CFT (Cardy--Peschel). There the free energy of a boundary with a corner acquires a term proportional to $c\log L$ with an angle-dependent coefficient. Here the same logarithm appears as the failure of the Hilbert--Schmidt condition, and the coefficient $\int_0^{\pi/2}\big(\sqrt\lambda/(\cos\vartheta+\lambda\sin\vartheta)-1/(\cos\vartheta+\sin\vartheta)\big)^2d\vartheta$ plays the role of the angle dependence. I have not compared the two coefficients. Together with (c), a corner is not removable by a local symmetry, but it is removable up to an invertible point-like insertion, namely the proper soliton $\sigma_\lambda$ of [DIT]. The corner is a solitonic, non-DHR excitation sitting at the corner and labelled by the velocity jump. So all the junction counts of this section (boundary conditions, T-junctions, BCC fields) are unchanged when the junction sits at a corner, while the corner itself carries a universal, boundary-condition-independent charge, by (a). The same applies to a kink in a transparent defect line, with $A$ replaced by $\cA$.

\emph{Status.} (a) is proved, including the vacuum-sector axiom (Lemma~\ref{lem:transport}). (b) is proved for every conformal net satisfying the axioms of [DIT], by reduction to their Theorem~3.6, and checked directly for the free fermion. (c) is proved: the sector axioms and invertibility follow from Lemma~\ref{lem:transport}. One standing assumption is used throughout: morphisms of bicolored intervals preserve the local coordinate at the colour point, as in [BDH] and Section~\nref{sec:net}.

\section{Lightlike boundaries}\label{sec:null}
[BKLR2] define a boundary to be a timelike line (\S1: ``a timelike plane of codimension~1 \dots\ such as $x=0$''). They observe that a transparent phase boundary may be placed on any timelike curve (\S3.6: ``the actual location of the boundary is perfectly arbitrary, and it may be any timelike curve'', the topological property). They do not discuss lightlike boundaries. Here is what the present framework gives. Use light-cone coordinates $u=t+x$, $v=t-x$, so that right-moving fields depend on $u$ and left-moving fields on $v$. Take the null line $\{v=0\}$, with $\cB^L$ on $\{v>0\}$ and $\cB^R$ on $\{v<0\}$.

\begin{proposition}[transparent null boundaries]\label{prop:null}
Let $\cB^L=\cA_2\rtimes Z[\bfA]$, $\cB^R=\cA_2\rtimes Z[\bfB]$, and let $D$ be the universal construction of [BKLR2] ($Z[\bfA]\times Z[\bfB]$).
\begin{enumerate}[nosep,label=(\alph*)]
\item The null boundary is consistent with the locality requirements of [BKLR2]. If $O_1\subset\{v>0\}$ and $O_2\subset\{v<0\}$ are spacelike separated, then $O_1$ lies to the spacelike left of $O_2$, so left-locality applies; all other pairs are timelike separated and impose nothing.
\item A double cone $O=I\times J$ ($u\in I$, $v\in J\ni0$) meeting the null line is cut by it into the two double cones $I\times J_\pm$, $J_\pm=J\cap\{\pm v>0\}$, and
\[
\cB^L(I\times J_+)\vee\cB^R(I\times J_-)=D(I\times J).
\]
So on regions meeting the boundary the null-boundary net is the universal net itself, as in the time-zero pinning (Proposition~\ref{prop:pinning}). The boundary conditions are its minimal central projections $E_{\mathbf m}$, the same as for timelike boundaries.
\item Consequently every statement above that is a statement about the universal nets and their centres and relative commutants carries over verbatim to null boundaries. This covers the classification by bimodules, the three-fold and $k$-fold products, the summands $D_{(\mathbf m_1,\mathbf m_2)}$ and $D_{(\mathbf m_1,\mathbf m',\mathbf m_2)}$, and the Y- and X-junction spaces.
\end{enumerate}
\end{proposition}
\begin{proof}
(a) Points $(u_1,v_1)$ with $v_1>0$ and $(u_2,v_2)$ with $v_2<0$ are spacelike separated iff $(u_1-u_2)(v_1-v_2)<0$, i.e.\ iff $u_1<u_2$, which says that the first is to the spacelike left of the second. (b) $\{v=0\}$ meets $I\times J$ in $I\times\{0\}$, and the two pieces are products of intervals. The inclusion $\subset$ is clear. For $\supset$: $\cA_2(I\times J)=A(I)\otimes A(J)=A(I)\otimes(A(J_+)\vee A(J_-))$ by strong additivity of $A$. Moreover $\cB^L(I\times J)$ is generated by $\cA_2(I\times J)$ and one charged intertwiner, which by transportability inside $I\times J$ may be chosen localized in $I\times J_+$; likewise for $\cB^R$ with $I\times J_-$. So $\cB^L(I\times J)\vee\cB^R(I\times J)=D(I\times J)$ is contained in the left side. (c) is immediate from (b).
\end{proof}
\emph{What does not carry over automatically.} The time-zero dictionary (Section~\ref{sec:apps}) must be replaced by a light-ray one. For fixed $I$, the assignment $J\mapsto D(I\times J)$ is a BDH-type defect on the $v$-light ray with the right-moving algebra $A(I)$ as an inert tensor factor. Identifying BDH fusion along the light ray with the universal construction would need Theorem~\nref{thm:A} for nets with such a spectator factor. The proof should go through unchanged, since $A(I)$ is inert, but I have not written it. Through the universal construction, juxtaposition of null boundaries is defined and gives the same results.

\textbf{Hard boundaries cannot be null.} In the Longo--Rehren construction a hard boundary is a timelike curve $u=\varphi(v)$, $\varphi\in\mathrm{Diff}_+$, along which the two chiral halves are glued ($T=\bar T$). On the null line $v=0$ the left-movers propagate \emph{along} the boundary and the right-movers cross it, so incoming right-moving energy cannot be reflected into left-movers. The gluing degenerates ($\varphi'\to0$ or $\infty$). In the language of Proposition~\ref{prop:corner}, bending a timelike hard boundary towards a null direction is a corner whose ratio $\lambda$ tends to $0$ or $\infty$. These values lie outside the family of solitons $\sigma_\lambda$, $\lambda\in\mathbb R_+$, which are pairwise inequivalent [DIT]. So a null hard boundary is a singular limit, not a boundary condition.

\textbf{Null interfaces between different chiral algebras.} A timelike interface is crossed by both chiralities, and transparency requires both chiral algebras to be common. A null interface $\{v=0\}$ is crossed only by the right-movers. So it must transmit $A_+$, but the left-moving nets on the two sides need not agree. Such an interface joins theories containing $A_+\otimes A^L_-$ and $A_+\otimes A^R_-$ with $A^L_-\ne A^R_-$, and in the left-moving direction it is a BDH defect between two \emph{different} chiral nets, tensored with the identity on $A_+$. This class is outside [BKLR2], where both chiral algebras are common, and outside Theorem~\nref{thm:A}, which fuses defects of a single net. It is the natural home of interfaces that preserve only one chiral algebra.

\section{Null interfaces between different chiral theories}\label{sec:nullint}
\textbf{Setting.} Fix a right-moving chiral net $A_+$ and two left-moving chiral nets $A^L_-$, $A^R_-$, all completely rational, with modular categories $\mathcal C_+$, $\mathcal C^L_-$, $\mathcal C^R_-$. Let $\cB^L\supset A_+\otimes A^L_-$ and $\cB^R\supset A_+\otimes A^R_-$ be local two-dimensional extensions. A \emph{null interface} along $\{v=0\}$ is a realization of $\cB^L$ on $\{v>0\}$ and $\cB^R$ on $\{v<0\}$ on a common Hilbert space, with the following properties. The two copies of $A_+$ coincide (the right-movers cross the interface and must be transmitted). Left-moving fields at different $v$ commute ($v$-locality across the interface). Left-locality holds for spacelike pairs as in Proposition~\ref{prop:null}(a). The interface algebra of a straddling double cone is $D(I\times J):=\cB^L(I\times J_+)\vee\cB^R(I\times J_-)$. Because the left-movers propagate along $\{v=0\}$ and never cross it, \emph{no relation between $A^L_-$ and $A^R_-$ is required}. This is what distinguishes null interfaces from timelike ones.

\begin{proposition}[Witt constraint]\label{prop:witt}
If $\cB^L$ and $\cB^R$ are maximal with chiral algebras exactly $A_+\otimes A^L_-$, resp.\ $A_+\otimes A^R_-$, then $[\mathcal C^L_-]=[\mathcal C_+]=[\mathcal C^R_-]$ in the Witt group of non-degenerate braided fusion categories. In particular the multiplicative central charges agree, $\xi(\mathcal C^L_-)=\xi(\mathcal C^R_-)$, and hence $c^L_-\equiv c^R_-\pmod 8$, using the standard relation $\xi(\mathrm{DHR}(A))=e^{2\pi ic/8}$ (quoted).
\end{proposition}
\begin{proof}
A maximal local extension of $A_+\otimes A_-$ with these chiral algebras is given by a Lagrangian $Q$-system in $\mathcal C_+\boxtimes\mathcal C_-^{\mathrm{rev}}$ [M\"uger], [BKL16]. By [DMNO], $\mathcal C\boxtimes\mathcal D^{\mathrm{rev}}$ has a Lagrangian algebra iff $[\mathcal C]=[\mathcal D]$. Apply this to both sides and use transitivity. $\xi$ is a Witt invariant [DMNO].
\end{proof}
So the constraint concerns the \emph{bulks}, not the interface: it allows left-moving algebras with different central charge (differing by multiples of $8$) and with no common subnet. That is the case not covered by anything in these notes.

\textbf{Case 1: a common left-moving subnet.} Suppose $A^L_-,A^R_-\supset C_-$, a common finite-index subnet, and suppose the interface is transparent to $A_+\otimes C_-$. Then $\cB^L,\cB^R$ are two local extensions of $\cA'_2:=A_+\otimes C_-$, and the null interface is a phase boundary relative to $\cA'_2$ in the general sense of [BKLR2] (\S3.6, where the common subnet need not be the full chiral algebra). The argument of Proposition~\ref{prop:null}(b) applies verbatim with $\cA'_2$ in place of $\cA_2$. Interfaces are classified by the minimal central projections of the braided product of the two $Q$-systems in $\mathrm{DHR}(\cA'_2)$, and they compose by the universal construction. For timelike lines, the time-zero net of $\cA'_2$ is $A_+\otimes\bar C_-$, Theorem~\nref{thm:A} applies to it as it stands, and Theorem~\nref{thm:fullc} does not (it uses $\mathcal C\boxtimes\mathcal C^{\rm opp}\simeq Z(\mathcal C)$). Computing the centres in this generality is a separate problem; for $Q$-systems that are themselves full centres relative to a common modular category, it reduces to Theorem~\nref{thm:kfold}. Example: a conformal embedding $C_-\subset A^R_-$ with $A^L_-=C_-$.

\textbf{Case 2 (outside the scope of these notes).} If $A^L_-$ and $A^R_-$ share no subnet, a null interface is a BDH defect between two different chiral nets with no common net in which it is a $Q$-system extension, and none of the methods here apply. We restrict to Case~1. Two remarks put Case~2 in context. First, defects between different nets exist and fuse associatively over finite-index middle nets [BDH, Conformal nets I, III, IV]. Second, Douglas--Henriques [DH] study \emph{CN-equivalence}: two nets are equivalent if there are defects between them whose composites are Morita equivalent to identity defects. They show that the representation category is a CN-invariant, and that a net is CN-invertible iff its $\mu$-index is $1$. They conjecture that the free fermion nets $\mathrm{Fer}(n)$ and $\mathrm{Fer}(n+576)$ are CN-equivalent, which would make the central charge \emph{not} a CN-invariant. They prove that if $\mathrm{Fer}(n)$ is CN-equivalent to the trivial net then $24\mid n$. So whether an interface between nets of different central charge can be invertible is a known question: conjecturally yes, in specific cases, subject to topological ($\mathrm{TMF}$-type) constraints. I have not found a more recent resolution.

\textbf{Decoupled interfaces.} For any two chiral nets $L,M$, the \emph{decoupled} assignment $D^{\rm dec}_{LM}(I):=L(I_\circ)\bar\otimes M(I_\bullet)$ (with $L$, $M$ on monochrome intervals) is a totally reflecting wall: nothing is transmitted. It satisfies isotony, locality and strong additivity, and its point algebra is trivial, $L(I_\circ)'\cap M(I_\bullet)'\cap D^{\rm dec}_{LM}(I)=\mathbb C$. I expect the vacuum-sector axiom to hold because the structure is the vacuum one near the fixed points of the reflection, but I have not checked it against the precise formulation of [BDH].

\begin{proposition}[fusion of decoupled interfaces]\label{prop:dec}
Assume $D^{\rm dec}_{LM}$ and $D^{\rm dec}_{MR}$ are defects. Then, on genuinely bicolored intervals,
\[
(D^{\rm dec}_{LM}\cfus_MD^{\rm dec}_{MR})(I)\ \cong\ L(I_\circ)\ \bar\otimes\ M(S)\ \bar\otimes\ R(I_\bullet),
\]
where $S$ is the strip between the two walls, $M(S)$ acts on the vacuum space of $M$, and the outer factors act on $L^2(L(I_\circ))$, $L^2(R(I_\bullet))$. Unless $M=\mathbb C$ this defect is not isomorphic to $D^{\rm dec}_{LR}$, and its point algebra is the type~III$_1$ factor $M(S)$.
\end{proposition}
\begin{proof}
Connes fusion over $M(J)$ acts only on the middle tensor factors, so with the charts of Definition~\nref{def:fusion}
\[
H_1\fus_{M(J)}H_2=L^2(L(I_\circ))\otimes\big(L^2(M(I^{++}_\bullet))\fus_{M(J)}L^2(M({}^{++}I_\circ))\big)\otimes L^2(R(I_\bullet)).
\]
By Haag duality inside an interval, $D^{\rm dec}_{LM}(I^{++})\cap M(\tilde J)'=L(I_\circ)\bar\otimes M(K_\bullet)$, and likewise for the other factor. The middle factor, with $M(K_\bullet)\vee M(L_\circ)$ acting on it, is exactly the situation of Proposition~\nref{prop:W} for the net $M$; its proof uses only $M(K)$, $M(\tilde J)$, $M(L)$ on vacuum spaces. So it is the vacuum space of $M$, with $M(K_\bullet)\vee M(L_\circ)=M(h(K_\bullet)\cup L_\circ)=:M(S)$. Now suppose $\varphi$ were a natural isomorphism onto $D^{\rm dec}_{LR}$. Naturality fixes $\varphi$ on $L(J)\otimes1\otimes1$ for white $J\subset I$ and on $1\otimes1\otimes R(J')$ for black $J'\subset I$, and these generate $L(I_\circ)\otimes1$ and $1\otimes R(I_\bullet)$. Hence $\varphi(1\otimes M(S)\otimes1)$ commutes with $L(I_\circ)\bar\otimes R(I_\bullet)$, a factor, so it is scalar. That contradicts injectivity unless $M(S)=\mathbb C$.
\end{proof}
So two totally reflecting walls cannot be fused away: the strip of $M$ between them survives as a type~III$_1$ point algebra, an ``$M$-bubble'', whatever the DHR category of $M$. This corrects the guess, made earlier in this work, that for holomorphic $M$ the fusion would be decoupled again. It contrasts with the transparent interfaces of Case~1, whose fusions are finite direct sums with finite-dimensional relative commutants. Decoupled interfaces are not transparent to any common subnet, so they lie outside the class to which we restrict.

\begin{remark}[a Haagerup example, conditional]\label{rem:haagerup}
A recent preprint [Hu] (not yet fully checked by a human reader) constructs a chiral algebra $U$ at $c=8$ with $\mathrm{Rep}(U)\simeq Z(\mathcal H_3)$, the Drinfeld centre of the Haagerup fusion category, and a commutative $Q$-system $A_3\in\mathrm{Rep}(U)$ whose extension is $(E_8)_1$, with $\mathrm{Rep}(U)_{A_3}\simeq\mathcal H_3$. If this holds, the results of these notes apply with $\mathcal C=\mathrm{Rep}(U)$ and give the following. The phases of the full theory over $U\otimes\bar U$ include the diagonal one ($Z[\mathbf 1]$) and the $E_8$ one ($Z[A_3]$). The transparent boundaries between them are labelled by $\mathrm{Irr}(\mathcal C_{A_3})=\mathrm{Irr}(\mathcal H_3)$, six boundary conditions. The $U\otimes\bar U$-preserving transparent defects of the $E_8$ phase are labelled by $\mathrm{Irr}({}_{A_3}\mathcal C_{A_3})=\mathrm{Irr}(\mathcal H_3\boxtimes\mathcal H_3^{\rm op})$. By Theorem~\nref{thm:main}, fusing two of them through the $E_8$ phase has relative commutant $\mathrm{End}(\mathbf m_1\otimes_{A_3}\mathbf m_2)$, governed by the Haagerup fusion rules. Since $\rho^2=\mathbf 1+\rho+\alpha\rho+\alpha^2\rho$ is multiplicity-free, these relative commutants are commutative. This is consistent with the claim of [Hu] that $\mathcal H_3$ acts on $(E_8)_1$ by topological defects, but it is not an independent check of that claim.
\end{remark}

\section{Hard boundaries meeting at a point}\label{sec:hardmeet}
For transparent configurations, every junction was treated by comparing two \emph{extended descriptions}: the configuration before the junction continued forward in time, and the one after it continued backward. Both agree with the physical theory on the causal complement of the junction. This works because the fields of a phase in $W_L\cup W_R$ are generated by $\cA_2$ and transportable charged fields, and $\cA_2$ passes through transparent boundaries unchanged. \emph{At a hard boundary it fails.} A chiral wave reflects there, so a field at a point remembers where its left-moving component was reflected. The extended future description reflects it at a fictitious continuation of the future boundary, and so disagrees with the physical theory even in $W_L$. (An earlier draft of this subsection used extended descriptions for hard boundaries; its proof of the splitting no-go compared the wrong representations, and is replaced below.) We therefore work with the physical configuration directly. Its chiral content is fixed by \emph{ray tracing}: follow each chiral component back along its light ray until it is either incoming from infinity or reflected at a hard boundary, where the gluing $T=\bar T$ of [LR04] turns it into the other chirality.

\textbf{Setup.} As in Section~\ref{sec:crossing}, $\ell,\ell'$ cross at $p=0$, giving regions $\Lambda_L,\Lambda_R,\Lambda_\pm$. The occupied set is $S$, hard boundaries separate occupied regions from empty ones, and transparent lines separate occupied regions from each other. A \emph{point-localized operator at $p$} is a bounded operator on the physical Hilbert space commuting with all fields on double cones in the occupied part of $W_L\cup W_R$. A hard boundary condition of the phase $\bfA$ is a left $\bfA$-module $\mathbf N$ in $\mathcal C$ ($\bfA$-modules are the boundary conditions of Theorem~\ref{thm:bcft}). Its chiral algebra is $Q_{\mathbf N}=\underline{\mathrm{End}}(\mathbf N)\in\mathcal C$, its boundary net is that of [LR04] on the vacuum space $H_{\mathbf N}$ of the chiral extension by $Q_{\mathbf N}$, and $F\colon\mathcal Z\to\mathcal C$, $\sigma\boxtimes\bar\tau\mapsto\sigma\otimes\tau$, is the unfolding (forgetful) functor. We quote:
\begin{itemize}[leftmargin=2.6em]
\item[(B)] [LR04], [LR09]. On a double cone $\{u\in I,\,v\in J\}$ of the half-system, unfolded with $I<J$ and with right-movers carried to the light ray by the gluing map, the boundary net is generated by $A(I)\vee A(J)$ and the charged fields of the extension by $F(Z[\bfA])$. On $H_{\mathbf N}$ these act through the bulk-to-boundary algebra map $\beta_{\mathbf N}\colon F(Z[\bfA])\to Q_{\mathbf N}$, by left multiplication on $Q_{\mathbf N}$.
\end{itemize}
By Proposition~\nref{prop:T8} (with $\mathbf X=\bfA$, $\mathbf Y=\mathbf 1$, $\mathcal B={}_{\bfA}\mathcal C$), $F(Z[\bfA])\cong\bigoplus_{\mathbf n}\underline{\mathrm{End}}(\mathbf n)$, a sum over the irreducible boundary conditions. The isomorphism is the action of the full centre on each $\mathbf n$, so $\beta_{\mathbf N}$ is the projection onto the summand $\mathbf n=\mathbf N$.

\begin{lemma}[causal complement of a boundary point]\label{lem:wedges}
Let $p$ lie on a hard boundary of a half-system. Ray tracing shows that the fields on double cones in the wedge $W$ spacelike to $p$ depend only on the past part of the boundary, so let $\mathbf N$ be the boundary condition there. Let $M_W$ be the von Neumann algebra these fields generate.
\begin{enumerate}[nosep,label=(\alph*)]
\item In every representation whose restriction to the chiral observables is a direct sum of DHR representations, $M_W$ contains $\pi(A(\mathbb R_-))\vee\pi(A(\mathbb R_+))$, whose commutant is $\pi(A)'$; in particular $M_W$ is of type~I.
\item In the representation corresponding to the module $\underline{\mathrm{Hom}}(\mathbf N',\mathbf N)$, which is the vacuum $H_{\mathbf N}$ for $\mathbf N'=\mathbf N$ and a boundary-condition-changing sector otherwise, $M_W'\cong\mathrm{End}_{\bfA}(\mathbf N')$. In particular $M_W$ is irreducible on $H_{\mathbf N}$ for simple $\mathbf N$.
\end{enumerate}
In contrast, a wedge $W$ in the bulk of a transparent configuration generates a type~III algebra.
\end{lemma}
\begin{proof}
\emph{Ray tracing.} Boost so that the past boundary is $x=0$, with the system on $x<0$, and let $q\in W=\{u<0<v\}$. Its right-moving component (constant $v>0$) comes from $x=-\infty$ and has not been reflected. Its left-moving component (constant $u<0$) was reflected at the boundary at time $u/(1+a)<0$, i.e.\ on the past part. Unfolding with the past gluing map therefore sends $W$ to $\mathbb R_-\cup\mathbb R_+$, the light ray with the single point $0$ removed, whatever the boundary does after $p$.

(a) By (B), $M_W\supset\pi(A(\mathbb R_-))\vee\pi(A(\mathbb R_+))$. For $\pi=\rho$ localized in $\mathbb R_+$, the commutant of this algebra is $A(\mathbb R_+)\cap\rho(A(\mathbb R_+))'=\mathrm{Hom}(\rho,\rho)$, by Haag duality and the standard fact that self-intertwiners of a localized endomorphism are determined on any interval containing its localization region. For a direct sum the commutant is $\pi(A)'$, which is of type~I, so $M_W$ is of type~I.

(b) The charged fields in $W$ are transported within $W$ by unitaries of $A(\mathbb R_\mp)\subset M_W$. Hence $M_W'$ consists of the operators in $\pi(A)'=\mathrm{End}_{\mathcal C}(\text{object})$ that commute with one charged field, i.e.\ by [BKLR, Ch.~4] the $F(Z[\bfA])$-module endomorphisms of $\beta_{\mathbf N}^*\underline{\mathrm{Hom}}(\mathbf N',\mathbf N)$. Since $\beta_{\mathbf N}$ is the projection onto a summand, these are the left $\underline{\mathrm{End}}(\mathbf N)$-module endomorphisms. By Ostrik's equivalence ${}_{\bfA}\mathcal C^{\rm op}\simeq\underline{\mathrm{End}}(\mathbf N)\text{-}\mathrm{Mod}$, $x\mapsto\underline{\mathrm{Hom}}(x,\mathbf N)$ (as in Theorem~\ref{thm:crossing}), they are $\mathrm{End}_{\bfA}(\mathbf N')$. For $\mathbf N'=\mathbf N$ this is $\mathrm{End}_{\underline{\mathrm{End}}(\mathbf N)}(\underline{\mathrm{End}}(\mathbf N))=\mathrm{Hom}(1,\underline{\mathrm{End}}(\mathbf N))=\mathrm{End}_{\bfA}(\mathbf N)$.

Bulk wedge: $\cA_2(W_L)=A(\mathbb R_-)\bar\otimes A(\mathbb R_+)$ is a type~III factor [GF]. It acts normally in DHR representations, and finite-index extensions with expectation of a type~III factor are of type~III [Kosaki], [Lo].
\end{proof}
Part (a) alone leaves the finite-dimensional algebra $\pi(A)'$, the chiral zero modes of the boundary. Part (b) shows that the bulk-to-boundary fields remove all of them except $\mathrm{End}(\mathbf N)$. This is the contrast with transparent defects. There, the analogous commutant is the relative commutant $\mathrm{End}(X)$ of Theorem~\nref{thm:kfold}, which carries the topological junction fields.

\begin{theorem}[hard boundaries meeting at a point]\label{thm:hardmeet}
Up to time reversal and reflection:
\begin{enumerate}[nosep,label=(\roman*)]
\item \emph{Points of a hard boundary, corners, and pinches.} Suppose $p$ lies on a hard boundary with simple condition $\mathbf N$ before $p$, and the configuration is realized in a representation of the form in Lemma~\ref{lem:wedges}(b) with simple $\mathbf N'$. This covers a straight boundary, a corner, and a change of boundary condition at $p$. Then the point-localized operators at $p$ are the scalars. For a \emph{pinch} ($S=\{L,R\}$), two half-systems with corners touching at $p$ realized on $H_L\otimes H_R$, the point-localized operators are $M_{W_L}'\bar\otimes M_{W_R}'=\mathbb C$. Nothing localized at a pinch can couple the two systems.
\item \emph{Splitting and joining} ($S=\{L,-,R\}$, resp.\ $\{L,+,R\}$). For splitting, ray tracing gives the left half-system after $p$ as follows: its chiral observables, unfolded, are $A_u$ on $\mathbb R_-$ and $A_v$ on $\mathbb R_+$, acting as two independent factors on the bulk Hilbert space. Hence, in any bulk state of vacuum type, the state of each half-system after the split is \emph{disjoint} from every positive-energy representation of that half-system: its chiral observables away from the split point generate the type~III algebra $A(\mathbb R_-)\bar\otimes A(\mathbb R_+)$, while in a positive-energy representation they generate a type~I algebra (Lemma~\ref{lem:wedges}(a)). For joining, the joined bulk state is disjoint from every bulk positive-energy representation, by the same comparison on a bulk wedge.
\item \emph{Creation, annihilation, passage} ($S\subset\{-,+\}$). The occupied part of $W_L\cup W_R$ is empty, so locality imposes nothing on the transition.
\item \emph{Corners with transparent decorations} (the remaining $S$, e.g.\ $\{L,-\}$, $\{L,-,+\}$). The physical theory is the corner theory of Proposition~\ref{prop:corner} with transparent lines ending on the boundary. Pushing the lines onto the boundary (Proposition~\ref{prop:T}, Theorem~\ref{thm:bcft}) gives a boundary condition $\mathbf m\otimes_{\bfB}\mathbf N$ before $p$, and (i) applies with that condition.
\end{enumerate}
\end{theorem}
\begin{proof}
(i) By Lemma~\ref{lem:wedges}(b) the commutant is $\mathrm{End}_{\bfA}(\mathbf N')=\mathbb C$. For the pinch, the two systems are independent by ray tracing (no chiral component crosses from one to the other), and the commutant of $M_{W_L}\otimes1\vee1\otimes M_{W_R}$ is $M_{W_L}'\bar\otimes M_{W_R}'$ by Tomita's commutation theorem.

(ii) Splitting: put the gap boundaries at $x=\mp at$, $t>0$, and let $q$ be in the left system after $p$. Its right-mover ($v>0$) is incoming from $x=-\infty$. Its left-mover with $u<0$ is incoming from the bulk past and never meets a boundary, since the line $u=\mathrm{const}<0$ does not meet $\{x=-at,\ t>0\}$. Its left-mover with $u>0$ is a right-mover with $v>0$ reflected at $x=-at$. So the left system's observables are generated by $A_u(\mathbb R_-)$ and $A_v(\mathbb R_+)$ of the bulk theory, acting on the bulk Hilbert space as independent tensor factors. They generate a normal image of the type~III factor $A(\mathbb R_-)\bar\otimes A(\mathbb R_+)$. A positive-energy representation of the half-system restricts on the unfolded chiral net to a sum of DHR representations, where the same observables generate a type~I algebra (Lemma~\ref{lem:wedges}(a)). Representations generating von Neumann algebras of different type on a common C$^*$-algebra are disjoint. Joining is the time reverse: after joining, a bulk wedge is generated by the observables of one incoming half-system, which are of type~I, whereas bulk positive-energy representations give type~III (Lemma~\ref{lem:wedges}).

(iii) $\Lambda_\pm\subset V_\pm$.

(iv) By ray tracing, the fields spacelike to $p$ see only the boundary condition before $p$.
\end{proof}

\begin{remark}[what (ii) does and does not localize; relation to local quenches]\label{rem:quench}
(a) Part (ii) is the operator-algebraic counterpart of the Calabrese--Cardy local quench [CC07], in which joining two half-lines needs a UV regulator $\epsilon$ and the joined state has infinite energy as $\epsilon\to0$. The disjointness is proved by comparing algebras of half-lines. On the unfolded light ray these touch \emph{both} at the split point $0$ and at $\infty$, so the proof does not by itself show where the singularity sits; Section~\ref{sec:smeared} separates the two.

(b) Localizing it at $p$ reduces to the statement that, for adjacent intervals $I_-=(-1,0)$, $I_+=(0,1)$, the product of the vacuum states on $A(I_-)$ and $A(I_+)$ does not extend to a normal state of $A(I_-\cup\{0\}\cup I_+)$, i.e.\ that the vacuum is not split across a point. For the free fermion this is proved in Section~\ref{sec:smeared} (Hilbert--Schmidt divergence). For general nets it is expected and not proved. Section~\ref{sec:smeared} also shows that smearing the junction over a bounded region removes the obstruction at the junction, but that on the infinite line an obstruction at spatial infinity remains.

(c) The pinch statement in (i) is a form of \emph{Haag duality at boundary points}: just as $(M_{W_L}\vee M_{W_R})'=\mathbb C$ for a bulk point, the causal complement of a boundary point leaves no room for point-localized operators. This is not automatic. The chiral observables alone would leave the finite-dimensional algebra $\pi(A)'$, and it is the bulk-to-boundary fields of (B) that remove it. A reducible boundary condition $\mathbf N=\bigoplus\mathbf N_i$ leaves exactly $\mathrm{End}_{\bfA}(\mathbf N)$, the projections onto its components.

(d) In particular, a change of boundary condition at a point is not implemented by any bounded point-localized operator. Boundary-condition-changing fields are unbounded fields of positive conformal weight, and the configuration with the change is a different sector, $\underline{\mathrm{Hom}}(\mathbf N',\mathbf N)$, not an operator on the vacuum sector. This agrees with Proposition~\ref{prop:T}, where topological boundary-condition-changing fields exist only between isomorphic conditions.
\end{remark}

\section{Smeared junctions of hard boundaries}\label{sec:smeared}
Theorem~\ref{thm:hardmeet}(ii) says that no state of vacuum type is carried by a hard-boundary split at a sharp point. Here a split is smeared over a region of size $a>0$: the fields far from the junction are required to be unaffected, and nothing is required inside the smeared region. We find that smearing does remove the type obstruction at the junction, but that on the infinite line an obstruction at spatial infinity remains. On a compact space, or on bounded regions, junctions exist and are classified, but they carry no topological data.

\textbf{Setting.} Work in the phase $\bfA$ with time-zero net $\cA_{\bfA}=\cA\rtimes Z[\bfA]$ (Section~\ref{sec:tz}, Lemma~\ref{lem:tz}), and let $H$ be its vacuum space; write $M(S):=\cA_{\bfA}(S)''\subset B(H)$ for an interval or half-line $S$ of the time-zero line, with its complement algebra $M(S)'=M(S^c)$ by Haag duality. By the time-slice property, $M(S)$ is the algebra of the causal completion of $S$. Let $K=H_{\mathbf N_-}\otimes H_{\mathbf N_+}$ be the vacuum space of two decoupled half-systems, the left one with hard boundary condition $\mathbf N_-$ on $(-\infty,-a_0]$ and the right one with $\mathbf N_+$ on $[a_0,\infty)$ (Section~\ref{sec:hard}), $0<a_0<a$. We quote:
\begin{itemize}[leftmargin=2.6em]
\item[(L)] [LR09], [LR04]. For an interval or half-line $S$ whose closure lies in the open region $x<-a_0$, the boundary net of the left system assigns to $S$ the algebra $\cA_{\bfA}(S)$, and the action of $\cA_{\bfA}(S)$ on $H_{\mathbf N_-}$ is faithful and normal. The same holds on the right.
\item[(S)] [DLR], [KLM]. If $A$ has finite maximal temperature (in particular if $A$ is completely rational), then for intervals $S_1,S_2$ with disjoint closures the pair $M(S_1),M(S_2)$ is a \emph{tensor pair}: the map $x_1x_2\mapsto x_1\otimes x_2$ extends to a normal isomorphism $M(S_1)\vee M(S_2)\to M(S_1)\bar\otimes M(S_2)$. The same holds for $\cA_{\bfA}$, a finite-index local extension.
\end{itemize}

\begin{definition}[smeared junction]\label{def:smeared}
Let $0<a<R\le\infty$ and $S_-=(-R,-a)$, $S_+=(a,R)$. An \emph{$(a,R)$-junction} is a unitary $V\colon H\to K$ with $V\,x\,V^*=\pi_-(x)$ for $x\in M(S_-)$ and $V\,y\,V^*=\pi_+(y)$ for $y\in M(S_+)$, where $\pi_\mp$ are the actions on the left, resp.\ right, factor of $K$ given by (L).
\end{definition}
The intertwining conditions say that, outside the smeared region $[-a,a]$ and inside $S_-\cup S_+$, the post-split state looks like the pre-split vacuum. Since the future wedge algebras are generated by the Cauchy data on $S_\mp$ at late times, and no characteristic from $S_\mp$ reaches the wall after it has formed, this is the condition that the fields there stream freely from the bulk past. The condition involves only $M(S_\mp)$, not $M(S_\mp)$ together with a neighbourhood of the wall, so the type obstruction of Lemma~\ref{lem:wedges} does not apply.

\begin{lemma}[reduction to independence]\label{lem:reduction}
An $(a,R)$-junction exists if and only if $(M(S_-),M(S_+))$ is a tensor pair in $B(H)$. In that case the junctions form a torsor: if $V$ is one, the others are $V u$, $u$ a unitary in $(M(S_-)\vee M(S_+))'=\cA_{\bfA}\big([-a,a]\big)''$ (for $R=\infty$; for finite $R$ the commutant is larger, $(M(S_-)\vee M(S_+))'=\cA_{\bfA}\big([-a,a]\cup\{|x|\ge R\}\big)''$).
\end{lemma}
\begin{proof}
($\Rightarrow$) $\pi_-\otimes\pi_+$ are commuting representations on the factors of $K$, so $M(S_-)\vee M(S_+)\cong\pi_-(M(S_-))\bar\otimes\pi_+(M(S_+))$ on $K$, with $x_-x_+\mapsto\pi_-(x_-)\pi_+(x_+)$; $V$ transports this to $H$. ($\Leftarrow$) $M(S_\mp)$ are type~III factors, so $M:=M(S_-)\bar\otimes M(S_+)$ is a type~III factor with separable predual. The representation $x_-x_+\mapsto x_-x_+$ on $H$ is a normal representation of $M$ by assumption, faithful since $M$ is a factor; $\pi_-\otimes\pi_+$ on $K$ is a normal faithful representation by (L). Normal faithful representations of a type~III factor on separable Hilbert spaces are unitarily equivalent. The torsor statement: $V,V'$ junctions give $V'V^*$ in the commutant of $\pi_-(M(S_-))\vee\pi_+(M(S_+))$, which $V^*$ transports to $(M(S_-)\vee M(S_+))'$ on $H$; the description of this commutant is Haag duality for $\cA_{\bfA}$ on the circle.
\end{proof}

\begin{theorem}[bounded smeared junctions]\label{thm:bounded}
Assume $A$ has finite maximal temperature. For $0<a<R<\infty$ there exist $(a,R)$-junctions for every pair $(\mathbf N_-,\mathbf N_+)$ of hard boundary conditions of $\bfA$, and they form a torsor over the unitary group of the type~III algebra $\cA_{\bfA}\big([-a,a]\cup\{|x|\ge R\}\big)''$.
\end{theorem}
\begin{proof}
$(a,R)$ with $R<\infty$ makes $S_-$, $S_+$ intervals with disjoint closures, so (S) and Lemma~\ref{lem:reduction} apply. (Intervals of the time-zero line, in the sense of the circle: $S_\pm$ stay away from the point at infinity.)
\end{proof}

\begin{proposition}[the closed string cut open]\label{prop:circle}
Let space be a circle of length $L$, with the chiral net $A$ on the circle and $H_{\rm circ}$ the vacuum space of $\cA_{\bfA}$. Cut the circle at $x=0$ into the interval $(0,L)$ with hard boundary conditions $\mathbf N_+$ at $0$ and $\mathbf N_-$ at $L$, and let $K_{\mathbf N_+\mathbf N_-}$ be the vacuum space of the resulting open system. For every $0<a<L/2$ there is a unitary $V\colon H_{\rm circ}\to K_{\mathbf N_+\mathbf N_-}$ intertwining the fields of the diamond $D_a$ over the arc $(a,L-a)$. The set of such $V$ is a torsor over the unitary group of $\cA_{\bfA}\big((-a,a)\big)''$, a type~III$_1$ factor.
\end{proposition}
\begin{proof}
Let $M_D:=\cA_{\bfA}\big((a,L-a)\big)''$. On $H_{\rm circ}$ it is a type~III$_1$ factor (algebra of an interval of the circle). By (L), the open system carries a faithful normal representation of the same abstract algebra, because the arc stays at distance $a$ from the ends. Both representations are of a type~III factor on separable spaces, hence unitarily equivalent. The commutant of $M_D$ on $H_{\rm circ}$ is $\cA_{\bfA}((-a,a))''$ by Haag duality, which gives the torsor.
\end{proof}
On the circle there is no infinity for the cut region to touch, so the junction exists without any condition on $a$. Both Theorem~\ref{thm:bounded} and Proposition~\ref{prop:circle} say the same thing: the condition that a junction be a unitary equivalence on the fields away from the junction is satisfied, in great abundance, by a torsor over the unitaries of the junction region. In particular \emph{it carries no topological data}. All the dependence on $(\mathbf N_-,\mathbf N_+)$ and on the phase is in which state $V\Omega$ is, in $K$, not in whether a $V$ exists.

\begin{proposition}[obstruction at infinity]\label{prop:infinity}
Let $R=\infty$. (a) An $(a,\infty)$-junction exists only if there is a normal state $\psi$ on $M(S_-)\vee M(S_+)\subset B(H)$ with $\psi(xy)=\langle\Omega,x\Omega\rangle\langle\Omega,y\Omega\rangle$ for $x\in M(S_-)$, $y\in M(S_+)$. (b) For the chiral free (Dirac) fermion, doubled to two dimensions, no such state exists, so there is no $(a,\infty)$-junction for any $a>0$.
\end{proposition}
\begin{proof}
(a) By Lemma~\ref{lem:reduction} the pair is a tensor pair, and then $\psi=\omega\otimes\omega$ is normal.

(b) The chiral right-moving part alone suffices. Let $I_-=(-\infty,-a)$, $I_+=(a,\infty)$, $I=I_-\cup I_+$. Let $P$ be the vacuum (Hardy) projection on $L^2(\mathbb R)$ with kernel $\frac1{2\pi i}\frac1{x-y-i0}$, and $C=\chi_IP\chi_I$, $C'=\chi_{I_-}P\chi_{I_-}+\chi_{I_+}P\chi_{I_+}$, acting on $L^2(I)$. The vacuum restricted to $\mathrm{CAR}(L^2(I))$ is the quasi-free state with two-point operator $C$. Any state $\psi$ with the product property of (a) is determined by it on products $xy$, hence is the quasi-free state with operator $C'$. If $\psi$ were normal, $\varphi_{C'}$ would lie in the folium of the factor state $\varphi_C$, hence be quasi-equivalent to it; by the criterion of Powers--St{\o}rmer and Araki [PS], $C^{1/2}-C'^{1/2}$ would then be Hilbert--Schmidt. Then so would be
\[
C-C'=C^{1/2}\big(C^{1/2}-C'^{1/2}\big)+\big(C^{1/2}-C'^{1/2}\big)C'^{1/2}.
\]
But $C-C'$ is the off-diagonal block with kernel $\frac1{2\pi i}\frac1{x-y}$, $x\in I_+$, $y\in I_-$ (and its adjoint), and
\[
\|C-C'\|_{\rm HS}^2=\frac2{4\pi^2}\int_a^\infty\!dx\int_{-\infty}^{-a}\!dy\ \frac1{(x-y)^2}=\frac1{2\pi^2}\int_a^\infty\frac{dx}{x+a}=\infty .
\]
\end{proof}
The divergence in (b) is logarithmic and sits at $x\to\infty$, where the two half-lines meet on the circle. The same computation for adjacent intervals $(-b,0)$, $(0,b)$ gives $\frac1{2\pi^2}\int_0^b\big(\frac1s-\frac1{s+b}\big)ds=\infty$, divergent at $s\to0$.

\begin{remark}[what smearing does and does not do]\label{rem:smeared}
(a) \emph{Where the singularity is.} The sharp-time obstruction of Theorem~\ref{thm:hardmeet}(ii) has two sources: the type mismatch at the junction point (type~III bulk wedge versus type~I wedge at a hard boundary), and the failure of independence at infinity (Proposition~\ref{prop:infinity}). Smearing removes the first (Lemma~\ref{lem:reduction}, Theorem~\ref{thm:bounded}, Proposition~\ref{prop:circle}) and cannot remove the second. The part of Remark~\ref{rem:quench}(a) saying that the singularity might sit at spatial infinity is thus correct: the exact statement ``the post-split state is the vacuum outside a bounded region'' is unattainable on the line. A physical quench state, for example $e^{-\epsilon H}$ applied to a cut state in [CC07], has finite energy and approaches the vacuum only asymptotically.

(b) \emph{Adjacent intervals.} The computation after Proposition~\ref{prop:infinity} proves for the free fermion the statement conjectured in Remark~\ref{rem:quench}(b): the vacuum is not split across a point. For general conformal nets I expect the same, but I have not proved it: it amounts to the statement that a half-sided modular inclusion of type~III$_1$ factors, such as $A(-\infty,-a)\subset A(-\infty,a)$, is never split. I have not found it in the literature.

(c) \emph{Energy.} The junctions of Theorem~\ref{thm:bounded} are all the unitaries of the junction region, so almost all of them have infinite energy for the future Hamiltonian, and nothing in this section selects a finite-energy one. Whether a finite-energy $V\Omega$ exists, and its dependence on $a$, is open. The expected answer, that the energy scales like $c/a$, is the CFT local-quench result [CC07], but I have not derived it here.

(d) \emph{Dictionary, not derived.} For a closed string cut open (Proposition~\ref{prop:circle}), the standard picture is that regularizing the entangling point by removing a disc of size $a$ and imposing boundary conditions $\mathbf N_\pm$ on it gives the open-string channel, with the entanglement spectrum organized by the sectors $i$ of multiplicity $n^{i}_{\mathbf N_-\mathbf N_+}=\dim\mathrm{Hom}_{\bfA}(\mathbf N_-\otimes i,\mathbf N_+)$ [OT15], [CT16]. In the present framework this would be a statement about the state $V\Omega\in K_{\mathbf N_+\mathbf N_-}$, i.e.\ about a particular element of the torsor, and not about the existence of $V$.
\end{remark}

\section{Regimes now covered}
Parallel transparent boundaries fused over $\cA$ or through an intermediate phase (Theorems~\nref{thm:main}, \ref{thm:through}); stacks of any length (Theorem~\nref{thm:kfold}); a transparent defect fused onto a parallel hard boundary (Theorem~\ref{thm:bcft}); Y-junctions (Proposition~\ref{prop:Y}, proved); T-junctions (Proposition~\ref{prop:T}, proved at the algebraic level, with the localization of the one-sided deformation not shown); X-junctions (Proposition~\ref{prop:X}: built from Y-junctions, proved; completeness via the sector calculus of Theorem~\ref{thm:sectors}; general junctions emit a defect $\mathbf a\in{}_{\bfA}\mathcal C_{\bfA}$). Intersecting transparent boundaries (Theorem~\ref{thm:crossing}): crossing fields characterized by causality as $\mathrm{Hom}_{\bfA\text{-}\bfC}(X,Y)$ in a fixed asymptotic channel, composition respected, channels matched across the crossing (Corollary~\ref{cor:channels}); horizontal composition at spacelike separation and the interchange law (Theorem~\ref{thm:horizontal}, Corollary~\ref{cor:interchange}); the structure is the 2-category of $Q$-systems and bimodules, modulo Lemmas~\ref{lem:unitins} and \ref{lem:unitins2}. Corners of hard boundaries and of defect lines (Proposition~\ref{prop:corner}): they factorize from all topological data and are equivalent but not isomorphic to no corner; the non-isomorphism is proved for all conformal nets (via [DIT]), and the corner carries the charge of a proper soliton labelled by the velocity jump. Lightlike transparent boundaries (Proposition~\ref{prop:null}): same classification and fusion results through the universal construction; the light-ray version of Theorem~\nref{thm:A} is not written. Null interfaces between different left-moving nets: Witt constraint (Proposition~\ref{prop:witt}); the case with a common subnet reduces to [BKLR2]; the case without one is left aside (Section~\ref{sec:nullint}). Hard boundaries meeting at a point (Theorem~\ref{thm:hardmeet}): no point-localized operators at boundary points, corners or pinches for simple boundary conditions (Lemma~\ref{lem:wedges}, via the bulk-to-boundary map); splitting or joining at a point produces a state disjoint from all positive-energy representations (localization at the split point reduces to non-splitness of the vacuum across a point, Remark~\ref{rem:quench}); creation and annihilation at a point are unconstrained by locality; decorated corners reduce to the corner theory. Smeared junctions (Section~\ref{sec:smeared}): the junction condition reduces to independence of the two outer half-lines (Lemma~\ref{lem:reduction}); on bounded regions and on a circle junctions exist for all boundary conditions and form torsors over the unitaries of the junction region, carrying no topological data (Theorem~\ref{thm:bounded}, Proposition~\ref{prop:circle}); on the infinite line an obstruction at spatial infinity remains (proved for the free fermion, Proposition~\ref{prop:infinity}); energy bounds are open. Not covered: interfaces that do not preserve both chiral algebras on timelike lines, quantitative corner data, and regulated (smeared) junctions of hard boundaries. Hard boundaries cannot be null.


\begin{thebibliography}{99}\small
\bibitem{BDH} A.~Bartels, C.~Douglas, A.~Henriques, \emph{Conformal nets III: fusion of defects}, arXiv:1310.8263; \emph{Conformal nets V: dualizability}, arXiv:1905.03393.
\bibitem{BKLR} M.~Bischoff, Y.~Kawahigashi, R.~Longo, K.-H.~Rehren, \emph{Tensor categories and endomorphisms of von Neumann algebras (with applications to quantum field theory)}, SpringerBriefs 2015, arXiv:1407.4793.
\bibitem{BKLR2} M.~Bischoff, Y.~Kawahigashi, R.~Longo, K.-H.~Rehren, \emph{Phase boundaries in algebraic conformal QFT}, Commun.\ Math.\ Phys.\ 342 (2016), arXiv:1405.7863.
\bibitem{BKL16} M.~Bischoff, Y.~Kawahigashi, R.~Longo, \emph{Characterization of 2D rational local conformal nets and its boundary conditions: the maximal case}, Doc.\ Math.\ 20 (2015); and with K.-H.~Rehren as above.
\bibitem{LR95} R.~Longo, K.-H.~Rehren, \emph{Nets of subfactors}, Rev.\ Math.\ Phys.\ 7 (1995).
\bibitem{LR04} R.~Longo, K.-H.~Rehren, \emph{Local fields in boundary conformal QFT}, Rev.\ Math.\ Phys.\ 16 (2004).
\bibitem{LR09} R.~Longo, K.-H.~Rehren, \emph{How to remove the boundary in CFT -- an operator algebraic procedure}, Commun.\ Math.\ Phys.\ 285 (2009).
\bibitem{FRS} J.~Fuchs, I.~Runkel, C.~Schweigert, \emph{TFT construction of RCFT correlators I}, Nucl.\ Phys.\ B 646 (2002).
\bibitem{Dav} A.~Davydov, \emph{Centre of an algebra}, arXiv:0908.1250.
\bibitem{DMNO} A.~Davydov, M.~M\"uger, D.~Nikshych, V.~Ostrik, \emph{The Witt group of non-degenerate braided fusion categories}, arXiv:1009.2117.
\bibitem{Ostrik} V.~Ostrik, \emph{Module categories, weak Hopf algebras and modular invariants}, Transform.\ Groups 8 (2003), arXiv:math/0111139.
\bibitem{EGNO} P.~Etingof, S.~Gelaki, D.~Nikshych, V.~Ostrik, \emph{Tensor categories}, AMS 2015.
\bibitem{KLM} Y.~Kawahigashi, R.~Longo, M.~M\"uger, \emph{Multi-interval subfactors and modularity of representations in conformal field theory}, CMP 219 (2001).
\bibitem{GF} F.~Gabbiani, J.~Fr\"ohlich, \emph{Operator algebras and conformal field theory}, CMP 155 (1993).
\bibitem{Kosaki} H.~Kosaki, \emph{Extension of Jones' theory on index to arbitrary factors}, J.\ Funct.\ Anal.\ 66 (1986).
\bibitem{CC07} P.~Calabrese, J.~Cardy, \emph{Entanglement and correlation functions following a local quench: a conformal field theory approach}, J.\ Stat.\ Mech.\ (2007) P10004.
\bibitem{Lo} R.~Longo, \emph{A duality for Hopf algebras and for subfactors I}, CMP 159 (1994).
\bibitem{CP} J.~Cardy, I.~Peschel, \emph{Finite-size dependence of the free energy in two-dimensional critical systems}, Nucl.\ Phys.\ B 300 (1988).
\bibitem{PS} R.~T.~Powers, E.~St{\o}rmer, \emph{Free states of the canonical anticommutation relations}, CMP 16 (1970); H.~Araki, \emph{On quasifree states of CAR and Bogoliubov automorphisms}, Publ.\ RIMS 6 (1970).
\bibitem{DIT} S.~Del Vecchio, S.~Iovieno, Y.~Tanimoto, \emph{Solitons and nonsmooth diffeomorphisms in conformal nets}, Commun.\ Math.\ Phys., doi:10.1007/s00220-019-03419-2, arXiv:1811.04501.
\bibitem{DH} C.~L.~Douglas, A.~Henriques, \emph{Topological modular forms and conformal nets}, in: Mathematical foundations of quantum field theory and perturbative string theory, Proc.\ Sympos.\ Pure Math.\ 83 (2011), arXiv:1103.4187.
\bibitem{Hu} T.-C.~Huang, \emph{Haagerup conformal field theory at central charge eight}, preprint, arXiv:2610.14441.
\bibitem{DLR} C.~D'Antoni, R.~Longo, F.~Radulescu, \emph{Conformal nets, maximal temperature and models from free probability}, J.\ Operator Theory 45 (2001).
\bibitem{OT15} Y.~Ohmori, Y.~Tachikawa, \emph{Physics at the entangling surface}, J.\ Stat.\ Mech.\ (2015) P04010, arXiv:1406.4167.
\bibitem{CT16} J.~Cardy, E.~Tonni, \emph{Entanglement Hamiltonians in two-dimensional conformal field theory}, J.\ Stat.\ Mech.\ (2016) 123103, arXiv:1608.01283.
\end{thebibliography}
\end{document}
